\documentclass[12pt,a4paper]{article}
\usepackage{ucebnice}

\title{Lineární rovnice, nerovnice a soustavy rovnic}
\author{Ondřej Škubala}
\date{\today}

\begin{document}

\maketitle
\newpage

\section{Úvod}

\begin{quotebox}
{„Možná se ti bude zdát téma této knihy obtížné, protože dosud není známé a začátečníkova mysl se snadno splete chybou; ale s tvou zvídavostí a mým vedením se ti jej podaří rychle ovládnout, neboť bystrý rozum podporovaný dobrým učením je nejkratší cestou k poznání.“}
{Diofantos z Alexandrie, \textit{Aritmetika (3. století n. l.)}}
\end{quotebox}

\noindent
Tato kapitola nás zavede do světa \important{rovnic a nerovnic} — jednoho z nejzákladnějších a zároveň nejmocnějších nástrojů celé algebry.  
Na příkladech lineárních rovnic si ukážeme, jak matematicky zachytit vztahy mezi neznámými veličinami, jak je analyzovat a jak hledat jejich řešení.  
Od jednoduchých rovnic o jedné neznámé přejdeme k soustavám rovnic, kde se neznámé prolínají navzájem, a nakonec i k rovnicím s absolutní hodnotou, které přinášejí do problému nový geometrický rozměr.  
Cílem není pouze „vyřešit“ rovnici, ale především \textbf{pochopit}, co samotný pojem rovnosti vyjadřuje a jaké důsledky má pro matematické myšlení.

\begin{historicalbox}
Pojem rovnice má své kořeny již ve starověku.  
V \textbf{egyptských papyrech} (např. v Ahmesově – Rhindově papyru, cca 1650 př. n. l.) nacházíme úlohy typu „něco a jeho sedmina dají dohromady 19“, které odpovídají dnešní lineární rovnici o jedné neznámé.  
Později \textbf{Babylonští matematici} dokázali řešit i kvadratické a soustavy rovnic pomocí tabulek a geometrických metod.  

Zásadní krok však přinesl \textbf{Diofantos z Alexandrie}, jehož dílo \textit{Aritmetika} (3. století n. l.) představuje první systematické pojednání o rovnicích a neznámých veličinách.  
Diofantos navázal na tradici řecké aritmetiky, ale překročil její geometrická omezení: začal místo geometrických představ pracovat se \important{symboly} a s abstraktním počtem.  
Zavedl označení pro neznámou, mocniny a některé operace, čímž vytvořil základy \textbf{symbolické algebry}.  
V jeho úlohách se často neřešily jen konkrétní čísla, ale hledaly se celé \textbf{množiny čísel, která splňují danou podmínku} — což dnes nazýváme \important{diofantické rovnice}.  
Tyto rovnice, v nichž se požaduje, aby řešení byla celá čísla, tvoří dodnes významnou oblast teorie čísel a přímé dědictví jeho práce.

O více než tisíc let později na jeho myšlenky navázal \textbf{Muhammad ibn Músá al-Chvárizmí}, který ve svém spisu \textit{Al-džabr wa’l-muqábala} popsal obecný postup pro řešení lineárních a kvadratických rovnic.  
Z jeho pojmu \textit{al-džabr} (doplňování, vyrovnávání) vznikl samotný název \textbf{algebra}.  
Teprve v 17. století \textbf{René Descartes} sjednotil geometrické a algebraické myšlení a zavedl symbolické zápisy rovnic, jak je známe dnes.  
Od té doby se algebra stala univerzálním jazykem, jímž dokážeme popsat prakticky jakýkoli vztah či jev ve světě.  

V českém prostředí se pojem \textbf{rovnice} objevuje až v raném novověku, kdy se začínají šířit překlady evropských učebnic matematiky.  
Zvláštní zásluhu na rozvoji matematického myšlení má \textbf{Jan Ámos Komenský}, jenž v \textit{Didaktice} zdůrazňuje význam názornosti a logické stavby učiva – zásady, které se naplno projevují právě v práci s rovnicemi.  
V 19. století pak čeští matematici jako \textbf{Bernard Bolzano}, \textbf{Karel Zahradník} či \textbf{František Josef Studnička} rozvíjejí přesnou logickou formulaci rovnic a vztahů mezi nimi, čímž dávají základ moderní výuce matematiky na našem území.  

Tak se příběh rovnic, započatý v prachu nilských polí a přepsaný v alexandrijských skriptoriích, uzavírá v prostředí našich školních lavic – a pokračuje v rukou každého, kdo se odhodlá porozumět tomu, co znamená, že dvě věci si jsou rovny.
\end{historicalbox}









\section{Lineární rovnice o jedné neznámé}


S lineárními rovnicemi o jedné neznámé jste se již setkali na základní škole, kde jste je řešili převážně mechanicky — přenášením členů a dělením koeficientem u neznámé.  
V této kapitole se však na rovnice podíváme \important{systematicky a matematicky}, tedy nejen s cílem „najít výsledek“, ale také pochopit, \textbf{proč} určité kroky děláme a jaké mají logické zdůvodnění.  
Zároveň se naučíme přesně formulovat, co znamená, že číslo je \textit{řešením} rovnice, a jak se tato řešení zapisují pomocí matematické symboliky.

\begin{definitionbox}
    \important{Lineární rovnice o jedné neznámé} je každá rovnice, kterou lze převést do tvaru
    \[
    ax + b = 0,
    \]
    kde $a, b \in \mathbb{R}$ a $a \neq 0$.
\end{definitionbox}


Lineární rovnice tohoto typu úzce souvisí s \textbf{lineární funkcí} $y = ax + b$, o které jsme psali v předchozí kapitole.  
Hledáme-li řešení rovnice $ax + b = 0$, ptáme se vlastně:  
\emph{„Pro jaké $x$ se hodnota funkce $y = ax + b$ rovná nule?“}  

Graficky to znamená, že hledáme \textbf{průsečík přímky} $y = ax + b$ \textbf{s osou $x$}, tedy bod, v němž přímka protíná osu $x$.  
Proto v rovnici symbolicky „dosazujeme“ $y = 0$.  
Řešení rovnice $ax + b = 0$ tak představuje \textbf{x-ovou souřadnici bodu, kde přímka protíná osu $x$}.  


Číslo $x$ označujeme jako \important{neznámou}, zatímco čísla $a$ a $b$ nazýváme \important{koeficienty} rovnice.  
Řešením rovnice rozumíme takové číslo $x_0$, které po dosazení do rovnice způsobí, že se obě její strany stanou \textbf{rovny}, tedy platí
\[
a x_0 + b = 0.
\]

Říkáme pak, že „\textbf{číslo $x_0$ rovnici splňuje}“ nebo také, že „\textbf{rovnici vyhovuje}“.  
Všechna čísla, která danou rovnici splňují, tvoří tzv. \important{množinu řešení} (též \textbf{množinu kořenů}) dané rovnice.  
Tuto množinu obvykle označujeme písmenem $\mathbf{M}$, případně $\mathbf{R}$ 
(z anglického \textit{roots}) či $\mathbf{K}$ (z českého \textit{kořeny}).  

\begin{examplebox}
    Například rovnice
    \[
    2x - 6 = 0
    \]
    má po úpravě $x = 3$.  
    Řekneme tedy, že číslo $3$ \textbf{rovnici splňuje}, a zapíšeme to jako
    \[
    M = \{3\}.
    \]
    Z geometrického hlediska to znamená, že přímka $y = 2x - 6$ protíná osu $x$ právě v bodě $(3, 0)$.

    \begin{center}
    \begin{axisgraph}[
        xmin=-1, xmax=6,
        ymin=-7, ymax=5,
        axis lines=middle,
        axis line style={thick},
        xlabel={$x$}, ylabel={$y$},
        title={Graf funkce $y = 2x - 6$},
        width=0.75\linewidth,
        height=6cm,
        samples=100,
        xtick={0,1,2,3,4,5,6},
        ytick={-6,-4,-2,0,2,4},
        grid=none
    ]
        % přímka
        \addplot[thick, blue, domain=-1:6] {2*x - 6};
        % průsečík
        \addplot[only marks, mark=*, blue] coordinates {(3,0)};
        \node[above right, blue] at (axis cs:3.2,0) {$[3;\,0]$};
    \end{axisgraph}
    \end{center}
\end{examplebox}

Vyřešit rovnici tedy znamená \important{najít všechna čísla, která tuto rovnici splňují}, a zapsat je jako množinu všech řešení.  
Tato množina může být prázdná, obsahovat jedno číslo nebo nekonečně mnoho čísel — podle toho, jaký typ rovnice řešíme.




\subsection{Ekvivalentní úpravy rovnic}

Při řešení rovnic se snažíme vyjádřit neznámou tak, aby stála „osamoceně“ na jedné straně rovnice.  
Abychom toho dosáhli, provádíme různé úpravy — přičítáme, odečítáme, násobíme či dělíme.  
Tyto úpravy jsou ale povoleny pouze tehdy, \important{pokud nezmění množinu řešení rovnice}.  
Každá úprava, která zachová rovnost obou stran, se nazývá \textbf{ekvivalentní úprava}.  


\begin{definitionbox}
    \textbf{Ekvivalentní úpravy rovnic jsou:}
    \begin{itemize}
        \item přičítání a odečítání reálného čísla,  
        \item přičítání a odečítání neznámé (nebo výrazu obsahujícího neznámou),  
        \item násobení reálným číslem kromě nuly,  
        \item násobení neznámou (nebo výrazem obsahujícím neznámou), pokud je různá od nuly (nebo výraz je různý od nuly),  
        \item dělení reálným číslem kromě nuly,  
        \item dělení neznámou (nebo výrazem obsahujícím neznámou), pokud je různá od nuly (nebo výraz je různý od nuly).
    \end{itemize}
\end{definitionbox}

Z matematického hlediska je každá z těchto úprav oprávněná, protože zachovává 
\textbf{rovnost} obou stran — tedy vztah, který říká, že levá a pravá strana mají stejnou hodnotu.  
Díky tomu se nemění ani množina všech čísel, která danou rovnici splňují.

\begin{examplebox}
    Vyřeš rovnici
    \[
    2x - 5 = \frac{x + 2}{2}.
    \]

    \textbf{Řešení:}

    \begin{alignat*}{3}
        2x - 5 &= \frac{x + 2}{2} &&\quad | \cdot 2\\
        2(2x - 5) &= 2 \cdot \frac{x + 2}{2}\\
        4x - 10 &= x + 2\\
        4x - 10 &= x + 2 &&\quad | -x\\
        4x - 10 - x &= x + 2 - x\\
        3x - 10 &= 2\\
        3x - 10 &= 2 &&\quad | +10\\
        3x - 10 + 10 &= 2 + 10\\
        3x &= 12 &&\quad | :3\\
        \frac{3x}{3} &= \frac{12}{3}\\
        x &= 4
    \end{alignat*}

    \textbf{Závěr:} Rovnice má jediné řešení $x = 4$, tedy množina řešení je
    \[
    K = \{4\}.
    \]
\end{examplebox}



Z geometrického hlediska platí, že každá ekvivalentní úprava \textbf{neposouvá průsečík přímky s osou $x$},  
ale pouze mění způsob, jakým je přímka algebraicky zapsána.  
Přičtení či odčítání členů tedy odpovídá posunu přímky nahoru nebo dolů,  
zatímco dělení či násobení mění její sklon, aniž by změnilo samotný bod průsečíku.

\begin{center}
    \begin{axisgraph}[
        xmin=-1, xmax=7,
        ymin=-8, ymax=8,
        axis lines=middle,
        xlabel={$x$}, ylabel={$y$},
        title={Rovnice $3x - 7 = 11$ a její ekvivalentní tvary},
        width=0.75\linewidth,
        height=6cm,
        samples=100,
        grid=none,
        xtick={0,2,4,6},
        ytick={-6,-3,0,3,6}
    ]
        % původní přímka
        \addplot[thick, gray!70!black, domain=-1:7] {3*x - 7};
        % posunutá (po +7)
        \addplot[dashed, gray!40, domain=-1:7] {3*x};
        % výsledná po dělení 3
        \addplot[thick, pastelorange!80!black, domain=-1:7] {3*x - 7 - 11};
        % průsečík
        \addplot[only marks, mark=*, pastelorange!80!black] coordinates {(6,0)};
        \node[above right, pastelorange!80!black] at (axis cs:2.2,0) {$[6;\,0]$};
    \end{axisgraph}
\end{center}

Každá z těchto úprav rovnice vede ke stejné odpovědi $x = 6$ — ať už rovnici zapisujeme 
jako $3x - 7 = 11$, nebo po úpravě $3x = 18$, nebo konečně $x = 6$.

\begin{notebox}
    Ekvivalentní úpravy fungují díky základní vlastnosti rovnosti:  
    \textit{jestliže jsou dvě veličiny stejné, zůstanou stejné i po přičtení, odčtení, násobení nebo dělení stejnou nenulovou hodnotou.}  
    Je to tedy obdoba principu zachování rovnováhy.
\end{notebox}



\subsection{Zkouška správnosti výsledku}

Vyřešit rovnici znamená najít všechna čísla, která ji splňují.  
Po nalezení řešení je proto vhodné provést tzv. \textbf{zkoušku}, 
při níž ověříme, že nalezené číslo skutečně vyhovuje původní rovnici.

Postupujeme tak, že dosadíme nalezené řešení (tzv. \textbf{kořen}) postupně do levé a 
pravé strany původní rovnice a vypočítáme jejich hodnoty:
\[
\begin{aligned}
L(\text{kořen}) &= [\text{dosazený výraz pro levou stranu}]\\
L(\text{kořen}) &= \text{výsledná hodnota levé strany}\\[6pt]
P(\text{kořen}) &= [\text{dosazený výraz pro pravou stranu}]\\
P(\text{kořen}) &= \text{výsledná hodnota pravé strany}\\[6pt]
L(\text{kořen}) &\stackrel{?}{=} P(\text{kořen})
\end{aligned}
\]


Pokud jsou obě hodnoty stejné, můžeme výsledek označit jako \textbf{správný}.  
Zapisujeme pak například:
\[
L(4) = P(4) = 3x - 7 = 11 \;\;\Rightarrow\;\; \text{rovnice je splněna.}
\]

\begin{examplebox}
Proveďte zkoušku pro rovnici
\[
2x - 5 = \frac{x + 2}{2}
\]
a nalezené řešení $x = 4$.

\textbf{Zkouška:}
\[
L(4) = 2 \cdot 4 - 5 = 8 - 5 = 3
\]
\[
P(4) = \frac{4 + 2}{2} = \frac{6}{2} = 3
\]
\[
L(4) = P(4)
\]

Proto je nalezené řešení $x = 4$ správné.
\end{examplebox}

\medskip
Zkouška není vždy nutná u jednoduchých rovnic, ale u složitějších výrazů
(zejména se zlomky, absolutními hodnotami nebo odmocninami) 
může pomoci odhalit chybu nebo \textbf{zdánlivé řešení}, 
které sice vychází výpočtem, ale původní rovnici ve skutečnosti nesplňuje.





\subsubsection*{Častá chyba: dělení výrazem, který může být roven nule}

Při řešení rovnic je nutné mít na paměti, že \textbf{dělení výrazem obsahujícím neznámou}
je povoleno pouze tehdy, pokud tento výraz \textbf{nemůže nabýt hodnoty nula}.  
V opačném případě bychom mohli některá řešení původní rovnice ztratit.

\medskip
Například u rovnice
\[
(x + 1)(x - 3) = x + 1
\]
se někteří studenti snaží rovnici zjednodušit dělením výrazem $(x + 1)$:
\[
(x + 1)(x - 3) = x + 1 \quad | : (x + 1)
\]
tím však dostanou pouze rovnici
\[
x - 3 = 1,
\]
jejímž řešením je $x = 4$.  
Jenže jsme zapomněli, že dělení $(x + 1)$ je povoleno pouze tehdy, pokud $x + 1 \neq 0$, tedy $x \neq -1$.  
Původní rovnice ale při dosazení $x = -1$ stále platí, a proto má ve skutečnosti \textbf{dvě řešení: $x = -1$ a $x = 4$}.  

\medskip
Správný postup, který zachovává všechna řešení, je například:

\begin{align*}
(x + 1)(x - 3) &= x + 1 &&\quad | - (x + 1)\\
(x + 1)(x - 3) - (x + 1) &= 0\\
(x + 1)\bigl((x - 3) - 1\bigr) &= 0\\
(x + 1)(x - 4) &= 0
\end{align*}

Tato rovnice má řešení $x = -1$ a $x = 4$, která odpovídají i původní rovnici.

\medskip
Abychom si ověřili, že obě tato řešení skutečně platí, provedeme \textbf{zkoušku správnosti:}

\textbf{Pro } $x = -1$:

\begin{align*}
L(-1) &= (-1 + 1)(-1 - 3) = 0 \cdot (-4) = 0\\
P(-1) &= (-1) + 1 = 0\\
L(-1) &\stackrel{?}{=} P(-1) \;\Rightarrow\; 0 = 0 \;\text{splňuje}
\end{align*}

\medskip
\textbf{Pro } $x = 4$:

\begin{align*}
L(4) &= (4 + 1)(4 - 3) = 5 \cdot 1 = 5\\
P(4) &= 4 + 1 = 5\\
L(4) &\stackrel{?}{=} P(4) \;\Rightarrow\; 5 = 5 \;\text{splňuje}
\end{align*}

\noindent
Vidíme, že obě nalezená řešení rovnici skutečně splňují.  
Pokud bychom ale rovnici upravili dělením $(x + 1)$, přišli bychom o první z nich,
protože bychom dělením odstranili situaci, kdy je $x = -1$.

\begin{notebox}
Při dělení výrazem obsahujícím neznámou vždy ověřte, zda tento výraz nemůže být roven nule.  
\textbf{Zkouška} je nejjistější způsob, jak ověřit, že nalezené řešení skutečně odpovídá původní rovnici.
\end{notebox}


\begin{examplebox}
Uvažujme rovnici se odmocninou:
\[
\sqrt{x + 16} = x + 4.
\]

\textbf{Řešení:}

Nejprve odstraníme odmocninu umocněním obou stran na druhou:
\begin{align*}
(\sqrt{x + 16})^2 &= (x + 4)^2\\
x + 16 &= x^2 + 8x + 16\\
0 &= x^2 + 7x\\
x(x + 7) &= 0
\end{align*}

\noindent
Odtud dostáváme dvě možná řešení:
\[
x_1 = 0, \quad x_2 = -7.
\]

\medskip
Protože jsme během řešení rovnici umocňovali, musíme provést \textbf{zkoušku}, abychom ověřili,
 zda obě hodnoty původní rovnici skutečně splňují.

\medskip
\textbf{Zkouška pro } $x = 0$:
\begin{align*}
L(0) &= \sqrt{0 + 16} = \sqrt{16} = 4\\
P(0) &= 0 + 4 = 4\\
L(0) &\stackrel{?}{=} P(0) \;\Rightarrow\; 4 = 4 \;\text{splňuje}
\end{align*}

\medskip
\textbf{Zkouška pro } $x = -7$:
\begin{align*}
L(-7) &= \sqrt{-7 + 16} = \sqrt{9} = 3\\
P(-7) &= -7 + 4 = -3\\
L(-7) &\stackrel{?}{=} P(-7) \;\Rightarrow\; 3 \neq -3 \;\text{nesplňuje}
\end{align*}

\noindent
Výsledek $x = -7$ tedy vznikl umocněním (které může rozšířit množinu řešení),  
a proto je pouze \textbf{zdánlivým řešením}.  
Správné řešení je pouze:
\[
K = \{0\}.
\]

\begin{notebox}
Při umocňování, odmocňování nebo jiných nelineárních úpravách může vzniknout tzv. \textbf{zdánlivé řešení}.  
Proto je vždy nutné provést zkoušku a ověřit, zda nalezené hodnoty původní rovnici opravdu splňují.
\end{notebox}
\end{examplebox}


\subsubsection*{Příklady řešení lineárních rovnic}

\begin{examplebox}
Vyřeš rovnici:
\[
2(x - 3) - (x - 4) = 3x - 10.
\]

\textbf{Řešení:}

\begin{align*}
2(x - 3) - (x - 4) &= 3x - 10\\
2x - 6 - x + 4 &= 3x - 10\\
x - 2 &= 3x - 10 &&| -x\\
-2 &= 2x - 10 &&| +10\\
8 &= 2x &&| :2\\
x &= 4
\end{align*}

\textbf{Zkouška:}
\begin{align*}
L(4) &= 2(4 - 3) - (4 - 4) = 2(1) - 0 = 2\\
P(4) &= 3 \cdot 4 - 10 = 12 - 10 = 2\\
L(4) &\stackrel{?}{=} P(4) \;\Rightarrow\; 2 = 2 \;\text{splňuje}
\end{align*}

Rovnice má jediné řešení $x = 4$.
\end{examplebox}


\begin{examplebox}
Vyřeš rovnici se zlomky:
\[
\frac{2x - 1}{3} = \frac{x + 2}{6}.
\]

\textbf{Řešení:}

Nejprve odstraníme zlomky — nejmenší společný jmenovatel je $6$.
\begin{align*}
\frac{2x - 1}{3} &= \frac{x + 2}{6} &&| \cdot 6\\
6 \cdot \frac{2x - 1}{3} &= 6 \cdot \frac{x + 2}{6}\\
2(2x - 1) &= x + 2\\
4x - 2 &= x + 2 &&| -x\\
3x - 2 &= 2 &&| +2\\
3x &= 4 &&| :3\\
x &= \frac{4}{3}
\end{align*}

\textbf{Zkouška:}
\begin{align*}
L\!\left(\tfrac{4}{3}\right) &= \frac{2 \cdot \tfrac{4}{3} - 1}{3} = \frac{\tfrac{8}{3} - 1}{3} = \frac{\tfrac{5}{3}}{3} = \tfrac{5}{9}\\
P\!\left(\tfrac{4}{3}\right) &= \frac{\tfrac{4}{3} + 2}{6} = \frac{\tfrac{4}{3} + \tfrac{6}{3}}{6} = \frac{\tfrac{10}{3}}{6} = \tfrac{5}{9}\\
L\!\left(\tfrac{4}{3}\right) &\stackrel{?}{=} P\!\left(\tfrac{4}{3}\right) \;\Rightarrow\; \tfrac{5}{9} = \tfrac{5}{9} \;\text{splňuje}
\end{align*}

Rovnice má řešení $x = \frac{4}{3}$.
\end{examplebox}


\begin{examplebox}
\textbf{Úloha:}  
Na parkovišti je celkem 50 vozidel.  
Počet automobilů je třikrát větší než počet motocyklů.  
Kolik je na parkovišti motocyklů a kolik automobilů?

\medskip
\textbf{Řešení:}

Označme počet motocyklů $x$.  
Potom počet automobilů je $3x$ a dohromady platí:
\[
x + 3x = 50.
\]
\begin{align*}
4x &= 50 &&| :4\\
x &= 12{,}5
\end{align*}

Protože počet vozidel musí být celé číslo, úloha v dané podobě nemá smysluplné řešení.  
Pokud by šlo o přibližný odhad, pak by na parkovišti bylo asi $13$ motocyklů a $37$ automobilů.

\begin{notebox}
Slovní úlohy často obsahují skrytou podmínku (např. že výsledek musí být celé nebo kladné číslo).  
Takové podmínky je třeba po vyřešení rovnice vždy zkontrolovat.
\end{notebox}
\end{examplebox}





\newpage
\section{Lineární nerovnice o jedné neznámé}


Zatímco rovnice vyjadřovala rovnost dvou výrazů, \textbf{nerovnice} popisuje situace, kdy je hodnota jednoho výrazu \emph{menší}, \emph{větší}, \emph{menší nebo rovna}, případně \emph{větší nebo rovna} hodnotě výrazu druhého.  
Výsledkem řešení nerovnice proto obvykle není jediné číslo, nýbrž \textbf{množina čísel}, která daný vztah splňují; v praxi ji zapisujeme intervalem a často ji také znázorňujeme na číselné ose či pomocí grafu související lineární funkce.

\begin{definitionbox}
\important{Lineární nerovnice o jedné neznámé} je každá nerovnice, kterou lze převést do tvaru
\[
ax + b \,\square\, 0,
\]
kde $a,b\in\mathbb{R}$, $a\neq 0$ a symbol $\square$ je jeden z $\{<,\,>,\,\le,\,\ge\}$.

Při úpravách používáme \textbf{ekvivalentní úpravy nerovnic}:
\begin{itemize}
  \item přičtení/odečtení stejného (reálného) čísla na obou stranách,
  \item přičtení/odečtení výrazu s neznámou na obou stranách,
  \item násobení/dělení obou stran \textbf{kladným} číslem (směr nerovnosti se \emph{nemění}),
  \item násobení/dělení obou stran \textbf{záporným} číslem (směr nerovnosti se \emph{obrací}),
  \item algebraické zjednodušování výrazů a odstraňování závorek.
\end{itemize}
\end{definitionbox}

\subsection*{Řešení krok za krokem a změna směru nerovnosti}
\begin{examplebox}
Vyřeš nerovnici se zlomkem:
\[
2x - 5 \le \frac{x + 2}{2}.
\]

\textbf{Řešení:}

Nejprve odstraníme zlomek – nejmenší společný jmenovatel je $2$.
\begin{align*}
2x - 5 &\le \frac{x + 2}{2} &&| \cdot 2\\
2(2x - 5) &\le x + 2\\
4x - 10 &\le x + 2 &&| -x\\
3x - 10 &\le 2 &&| +10\\
3x &\le 12 &&| :3\\
x &\le 4
\end{align*}

\noindent
\textbf{Výsledek:} $K = (-\infty,\,4\rangle$.
\end{examplebox}

\begin{examplebox}
Vyřeš nerovnici s neznámou na obou stranách:
\[
3x - 7 \le 5x - 13.
\]

\textbf{Řešení:}
\begin{align*}
3x - 7 &\le 5x - 13 &&| -3x\\
-7 &\le 2x - 13 &&| +13\\
6 &\le 2x &&| :2\\
3 &\le x
\end{align*}

\noindent
\textbf{Výsledek:} $K = \langle3,\infty)$.
\end{examplebox}

\subsubsection*{Typické chyby}
\begin{itemize}
  \item Zapomenutí \textbf{obrátit nerovnost} při dělení nebo násobení \textbf{záporným} číslem.
  \item Mechanické „přenášení členů“ bez ohledu na znaménko a následná záměna směru.
  \item Nesprávný zápis výsledku: chybné použití \emph{otevřeného} vs. \emph{uzavřeného} krajního bodu.
  \item U zlomků dělení výrazem, který může být nulový (v takovém případě je třeba pracovat s podmínkami).
\end{itemize}

\subsubsection*{Příklady řešení lineárních rovnic}

\begin{examplebox}
Vyřeš nerovnici
\[
5x - 7 \ge 3.
\]

\textbf{Řešení:}
\begin{align*}
5x - 7 &\ge 3 &&| +7\\
5x &\ge 10 &&| :5\\
x &\ge 2
\end{align*}

\noindent
\textbf{Výsledek:} $M = [2,\infty)$.
\end{examplebox}

\newpage
\begin{examplebox}
Vyřeš nerovnici
\[
2(x-3) - (x+1) < 5 - x.
\]

\textbf{Řešení:}
\begin{align*}
2x - 6 - x - 1 &< 5 - x\\
x - 7 &< 5 - x &&| +x\\
2x - 7 &< 5 &&| +7\\
2x &< 12 &&| :2\\
x &< 6
\end{align*}

\noindent
\textbf{Výsledek:} $M = (-\infty,6)$.
\end{examplebox}

\begin{examplebox}
Vyřeš nerovnici
\[
\frac{2x-1}{3} \le \frac{x+5}{6}.
\]

\textbf{Řešení:} Nejmenší společný jmenovatel je $6$.
\begin{align*}
\frac{2x-1}{3} &\le \frac{x+5}{6} &&| \cdot 6\\
6\cdot \frac{2x-1}{3} &\le 6\cdot \frac{x+5}{6}\\
2(2x-1) &\le x+5\\
4x - 2 &\le x + 5 &&| -x\\
3x - 2 &\le 5 &&| +2\\
3x &\le 7 &&| :3\\
x &\le \frac{7}{3}
\end{align*}

\noindent
\textbf{Výsledek:} $M = (-\infty,\,7/3]$.
\end{examplebox}

\begin{examplebox}
Vyřeš nerovnici
\[
-4x + 3 > 2x - 9.
\]

\textbf{Řešení:}
\begin{align*}
-4x - 2x &> -9 - 3 &&\quad \text{(převedeme vše s $x$ vlevo, čísla vpravo)}\\
-6x &> -12 &&| :(-6) \quad \text{(dělíme záporným číslem)}\\
x &< 2
\end{align*}

\noindent
\textbf{Výsledek:} $M = (-\infty,2)$.
\end{examplebox}

\begin{notebox}
\textbf{Proč se nerovnost při dělení záporným číslem obrací?}

Představme si dvě čísla $a < b$.  
Pokud od obou stran odečteme $a + b$, dostaneme
\[
a - a - b < b - a - b \;\Rightarrow\; -b < -a.
\]
Tím vidíme, že po \textbf{násobení nebo dělení záporným číslem} se směr nerovnosti skutečně \textbf{obrací}.  
To odpovídá i intuici: pokud čísla vynásobíme $-1$, zrcadlí se na číselné ose — menší číslo se stane větším a naopak.
\end{notebox}


\subsection*{Zápis a grafické znázornění výsledků}

Řešení nerovnic představuje \textbf{množinu čísel}, která splňují daný vztah.  
Tyto množiny zapisujeme pomocí \textbf{intervalů} a znázorňujeme na \textbf{číselné ose}.  

\medskip
\noindent
\textbf{Základní význam symbolů:}

\begin{center}
\renewcommand{\arraystretch}{1.4}
\begin{tabular}{ccl}
\textbf{Symbol} & \textbf{Význam} & \textbf{Zápis intervalu (evropský)}\\
$a < x < b$ & $x$ je mezi $a$ a $b$ & $\langle a;\,b\rangle$ — oba body otevřené\\
$a \le x \le b$ & $x$ je mezi $a$ a $b$, včetně krajů & $\langle a;\,b\rangle$ — oba body uzavřené\\
$x > a$ & $x$ je větší než $a$ & $\langle a;\,+\infty)$\\
$x \ge a$ & $x$ je větší nebo rovno $a$ & $\langle a;\,+\infty)$ (uzavřeno vlevo)\\
$x < a$ & $x$ je menší než $a$ & $(-\infty;\,a\rangle$\\
$x \le a$ & $x$ je menší nebo rovno $a$ & $(-\infty;\,a\rangle$ (uzavřeno vpravo)\\
\end{tabular}
\end{center}

\medskip
\noindent
\textbf{Značení na číselné ose:}

\begin{center}
\begin{tabular}{cc}
% --- 1. graf: x < 2 ---
\begin{axisgraph}[
    xmin=-3, xmax=7,
    ymin=0, ymax=1,
    axis x line=middle,
    axis y line=none,
    width=0.45\linewidth,
    height=2.2cm,
    ytick=\empty,
    xtick={-2,0,2,4,6},
    xticklabels={},
    clip=false
]
% interval vlevo
\draw[thick, pastelorange, <-] (axis cs:-2,0.5) -- (axis cs:1.87,0.5);
% otevřený bod
\draw[thick, pastelorange] (axis cs:2,0.5) circle (2.5pt);
% popisek
\node[above, pastelorange!80!black] at (axis cs:0,0.55) {$x<2$};
\end{axisgraph}
&
% --- 2. graf: x ≤ 2 ---
\begin{axisgraph}[
    xmin=-3, xmax=7,
    ymin=0, ymax=1,
    axis x line=middle,
    axis y line=none,
    width=0.45\linewidth,
    height=2.2cm,
    ytick=\empty,
    xtick={-2,0,2,4,6},
    xticklabels={},
    clip=false
]
\draw[thick, pastelblue, <-] (axis cs:-2,0.5) -- (axis cs:2,0.5);
\filldraw[pastelblue] (axis cs:2,0.5) circle (2.5pt);
\node[above, pastelblue!70!black] at (axis cs:0,0.55) {$x\le2$};
\end{axisgraph}
\\[1em]
% --- 3. graf: x > 4 ---
\begin{axisgraph}[
    xmin=-3, xmax=7,
    ymin=0, ymax=1,
    axis x line=middle,
    axis y line=none,
    width=0.45\linewidth,
    height=2.2cm,
    ytick=\empty,
    xtick={-2,0,2,4,6},
    xticklabels={},
    clip=false
]
\draw[thick, pastelgreen, ->] (axis cs:4.15,0.5) -- (axis cs:6,0.5);
\draw[thick, pastelgreen] (axis cs:4,0.5) circle (2.5pt);
\node[above, pastelgreen!80!black] at (axis cs:5,0.55) {$x>4$};
\end{axisgraph}
&
% --- 4. graf: x ≥ 4 ---
\begin{axisgraph}[
    xmin=-3, xmax=7,
    ymin=0, ymax=1,
    axis x line=middle,
    axis y line=none,
    width=0.45\linewidth,
    height=2.2cm,
    ytick=\empty,
    xtick={-2,0,2,4,6},
    xticklabels={},
    clip=false
]
\draw[thick, pastelviolet, ->] (axis cs:4,0.5) -- (axis cs:6,0.5);
\filldraw[pastelviolet] (axis cs:4,0.5) circle (2.5pt);
\node[above, pastelviolet!70!black] at (axis cs:5,0.55) {$x\ge4$};
\end{axisgraph}
\end{tabular}
\end{center}


\medskip
\noindent
\textbf{Interpretace:}
\begin{itemize}
  \item Otevřený bod (prázdné kolečko) znamená, že \textbf{krajní hodnota do množiny nepatří}.
  \item Plný bod znamená, že \textbf{krajní hodnota je součástí množiny řešení}.
  \item Směr šipky určuje, zda jsou řešením čísla \textbf{menší} (vlevo) nebo \textbf{větší} (vpravo) než daná hodnota.
  \item Dvojitá nerovnice (např. $1 < x \le 4$) odpovídá úseku $\langle 1;4\rangle$ — uzavřenému vpravo, otevřenému vlevo.
\end{itemize}

\begin{notebox}
\textbf{Poznámka k evropskému zápisu:}  
V české matematické praxi zapisujeme intervaly pomocí \textbf{hranatých} a \textbf{špičatých závorek}:  
$\langle a;b\rangle$, $\langle a;b)$, $(a;b\rangle$, $(a;b)$.  
Naopak souřadnice bodů zapisujeme jako \textbf{[x;y]} — hranaté závorky tedy patří bodům,  
špičaté intervalům.
\end{notebox}




\section{Grafické řešení lineárních rovnic a nerovnic}


Každou lineární rovnici nebo nerovnici můžeme chápat také \textbf{graficky}.  
Pokud si uvědomíme, že výrazy na obou stranách rovnice představují \textbf{funkce proměnné} \(x\),  
pak řešení rovnice odpovídá \textbf{průsečíku jejich grafů}.  

\medskip

Například rovnici
\[
2x + 1 = -x + 3
\]
můžeme chápat jako soustavu dvou lineárních funkcí:
\[
y_1 = 2x + 1, \qquad y_2 = -x + 3.
\]

\begin{examplebox}
\textbf{Řešení:}

Zapsaná rovnice \(2x + 1 = -x + 3\) říká, že hledáme takové \(x\), pro které jsou hodnoty obou funkcí stejné,  
tedy jejich grafy se protínají.

\begin{align*}
2x + 1 &= -x + 3\\
3x &= 2\\
x &= \tfrac{2}{3}
\end{align*}


\textbf{Řešením rovnice je tedy průsečík:}

\begin{center}
\begin{axisgraph}[
    xmin=-2, xmax=3,
    ymin=-1, ymax=5,
    axis lines=middle,
    axis line style={thick},
    xlabel={$x$}, ylabel={$y$},
    grid=none,
    width=0.75\linewidth,
    height=6.5cm,
    samples=100,
    domain=-2:3,
    legend style={at={(0.95,0.95)}, anchor=north east, font=\small}
]
% přímky
\addplot[thick, pastelblue, domain=-2:3]{2*x + 1};
\addplot[thick, pastelorange, domain=-2:3]{-x + 3};

% svislá pomocná čára k ose x
\addplot[dashed, gray!70] coordinates {(0.6667, 0) (0.6667, 2.3333)};

% průsečík
\addplot[only marks, mark=*, mark size=2.5pt, pastelgreen!70!black] coordinates {(0.6667, 2.3333)};

% popisek pouze na ose x
\node[below, gray!70!black] at (axis cs:0.6667, 0) {$\tfrac{2}{3}$};
\end{axisgraph}
\end{center}

\textbf{Interpretace:}  
Grafické řešení rovnice znamená najít bod, v němž se obě přímky protínají.  
Každá další rovnice představuje „posunutí“ nebo „naklonění“ těchto přímek –  
pokud mají různé směrnice, protínají se právě v jednom bodě; pokud jsou rovnoběžné, nemají řešení;  
a pokud splývají, mají nekonečně mnoho společných bodů, tedy nekonečně mnoho řešení.
\end{examplebox}


\begin{notebox}
\textbf{Shrnutí:}  
- Rovnice \(f(x) = g(x)\) má řešení v těch \(x\), kde se grafy funkcí \(y=f(x)\) a \(y=g(x)\) protínají.  
- Počet řešení závisí na poloze přímek:
  \begin{itemize}
    \item různé směrnice → jeden průsečík (jedno řešení),
    \item shodné směrnice, různé posunutí → žádný průsečík (žádné řešení),
    \item shodné přímky → nekonečně mnoho řešení.
  \end{itemize}
\end{notebox}


\newpage
\section{Lineárně lomené rovnice}

V této kapitole se zaměříme na rovnice, ve kterých se neznámá nachází ve \textbf{jmenovateli zlomku}.  
Takové rovnice označujeme jako \important{lineárně lomené} nebo též \important{rovnice v podílovém tvaru}.  

\medskip
Na první pohled vypadají složitěji, ale ve skutečnosti se po vhodných úpravách dají převést na běžné lineární rovnice.  
Jediným rozdílem je, že je nutné nejprve určit tzv. \textbf{podmínky řešitelnosti} — hodnoty, pro které \textbf{jmenovatel nesmí být roven nule}.

\begin{definitionbox}
\textbf{Lineárně lomená rovnice} je rovnice, která obsahuje zlomky s neznámou ve jmenovateli, např.:
\[
\frac{2x - 5}{x - 3} = 1.
\]
\end{definitionbox}

\textbf{Postup řešení:}
%\begin{enumerate}[label=\arabic*.]
%  \item Určíme \textbf{podmínky řešitelnosti}: výrazy ve jmenovateli nesmí být nulové.
%  \item Násobíme obě strany rovnice \textbf{nejmenším společným jmenovatelem (NSJ)}.
%  \item Získanou rovnici upravíme na běžnou lineární rovnici.
%  \item Vyřešíme ji a ověříme, zda získané řešení splňuje podmínky (nevede k dělení nulou).
%\end{enumerate}


\medskip 
Uvažujme rovnici
\[
\frac{2x - 5}{x - 3} = 1.
\]

Než začneme rovnici upravovat, stanovíme \textbf{podmínku řešitelnosti}.  
Ve jmenovateli se nachází výraz \(x - 3\), který nesmí být nulový:
\begin{align*}
  x - 3 &\neq 0\\
  x &\neq 3
\end{align*}

Tím jsme určili hodnotu, kterou \(x\) nesmí nabývat — jinak by došlo k dělení nulou.

\medskip
Nyní odstraníme jmenovatel.  
Rovnici vynásobíme výrazem \(x - 3\), který je (dle podmínky) nenulový:
\[
(x - 3)\cdot\frac{2x - 5}{x - 3} = (x - 3)\cdot 1.
\]

Po úpravě se jmenovatelé zkrátí:
\[
2x - 5 = x - 3.
\]

Získali jsme běžnou lineární rovnici, kterou už vyřešíme obvyklým způsobem:
\begin{align*}
  2x - 5 &= x - 3 \\
  2x - x &= -3 + 5 \\
  x &= 2
\end{align*}

Nakonec ověříme, že nalezené řešení neporušuje podmínku \(x \neq 3\).  
Protože \(x = 2\) je povolené, můžeme výsledek přijmout.

\[
\textbf{Řešení:} \quad x = 2.
\]

\[
\textbf{Množina řešení:} \quad K = \{2\}.
\]

\medskip
Vidíme, že kromě běžného postupu přibylo jen jedno pravidlo:  
vždy je třeba zkontrolovat, zda nalezené řešení nevede k dělení nulou.  
Pokud by tomu tak bylo, musíme jej \textbf{vyloučit} jako \textbf{zdánlivé řešení}.

\medskip
Podívejme se nyní na situaci, kdy nalezené řešení nevyhovuje podmínce — vznikne tak \textbf{zdánlivé řešení}.

\begin{examplebox}
\textbf{Příklad:}  
Vyřešte rovnici
\[
\frac{x + 2}{x - 1} = \frac{3x - 1}{x - 1}.
\]

\textbf{Podmínka:}  
Jmenovatel nesmí být nulový, tedy \(x - 1 \neq 0 \Rightarrow x \neq 1.\)

\textit{Rovnici vynásobíme jmenovatelem \((x - 1)\) a upravíme:}
\begin{align*}
(x - 1)\cdot \frac{x + 2}{x - 1} &= (x - 1)\cdot \frac{3x - 1}{x - 1}\\[-2pt]
x + 2 &= 3x - 1\\
3 &= 2x\\
x &= \tfrac{3}{2}.
\end{align*}

\textbf{Zkouška:}  
Pro \(x = \tfrac{3}{2}\) je jmenovatel \(x - 1 = \tfrac{1}{2} \neq 0\), takže řešení je platné.  
Pokud by vyšlo \(x = 1\), nebylo by přípustné (dělení nulou).

\[
\textbf{Řešení:}\; x = \tfrac{3}{2} \qquad
\textbf{Množina řešení:}\; K = \left\{\tfrac{3}{2}\right\}.
\]
\end{examplebox}


\begin{notebox}
\textbf{Shrnutí postupu:}
\begin{enumerate}
  \item Urč podmínky, za kterých mají všechny jmenovatele smysl (nesmí být nulové).
  \item Rovnici vynásob nejmenším společným jmenovatelem (NSJ) – tím odstraníš zlomky.
  \item Vyřeš vzniklou lineární rovnici.
  \item Každé nalezené řešení ověř: pokud ruší jmenovatel, je \textbf{zdánlivé}.
\end{enumerate}

\textbf{Pozor:} Každá rovnice s neznámou ve jmenovateli musí mít \textbf{zkoušku},  
jinak hrozí, že mezi řešeními zůstane hodnota, pro kterou je rovnice nedefinovaná.
\end{notebox}



\newpage
\section{Soustavy lineárních rovnic}
\subsection{Soustava dvou rovnic o dvou neznámých}


Zatímco při řešení jedné rovnice hledáme hodnotu jedné neznámé,  
v soustavách rovnic se snažíme nalézt \important{společné řešení} pro více neznámých najednou.  
Takové úlohy se v praxi objevují velmi často – od hledání průsečíku dvou přímek, přes rovnováhu sil v mechanice,  
až po analýzu hospodářských vztahů v ekonomii.  

\medskip
Každou z rovnic můžeme chápat jako \textbf{podmínku}, kterou mají neznámé splňovat.  
Hledáme tedy takové hodnoty, které vyhovují \textbf{všem} rovnicím zároveň.

\begin{definitionbox}
\textbf{Soustava dvou lineárních rovnic o dvou neznámých} je dvojice rovnic tvaru
\[
\begin{cases}
a_1x + b_1y = c_1,\\[4pt]
a_2x + b_2y = c_2,
\end{cases}
\]
kde \(a_1, b_1, c_1, a_2, b_2, c_2 \in \mathbb{R}\), přičemž alespoň jeden z koeficientů \(a_i, b_i\) je nenulový.

\medskip
\textbf{Řešením soustavy} nazýváme každou uspořádanou dvojici \([x_0;y_0]\), která splňuje \textbf{obě rovnice zároveň}.  
Množinu všech řešení označujeme písmenem \(K\).
\end{definitionbox}

\noindent
Je-li některá z rovnic soustavy neplatná nebo nadbytečná, může to ovlivnit existenci řešení celé soustavy:

\begin{itemize}
  \item Pokud \(a_1 = b_1 = c_1 = 0\), pak první rovnice neobsahuje žádnou informaci  
  a soustavu můžeme \textbf{zredukovat na jedinou rovnici} (tedy řešení přechází na soustavu s jednou rovnicí o dvou neznámých).  
  Takové případy budeme řešit dále při studiu \textbf{jedné lineární rovnice s více neznámými}.
  
  \item Pokud \(a_1 = b_1 = 0\) a zároveň \(c_1 \neq 0\), pak první rovnice představuje \textbf{nesmyslné tvrzení}  
  (například \(0 = 5\)), čili rovnice nemá řešení a tedy ani soustava \textbf{nemá žádné řešení}, protože nelze splnit všechny podmínky zároveň.
\end{itemize}


Každou z těchto rovnic můžeme chápat jako přímku v rovině.  
Řešení soustavy pak odpovídá \textbf{průsečíku těchto přímek} – tedy bodu, který \textbf{leží na obou přímkách zároveň}.

\subsubsection*{Z geometrického hlediska tedy platí:}

\begin{notebox}
Každá lineární rovnice o dvou neznámých představuje v rovině přímku.  
Řešení soustavy dvou lineárních rovnic o dvou neznámých odpovídá bodu (nebo více bodům), které leží \textbf{na obou přímkách zároveň}.
\end{notebox}

\begin{center}
\begin{tikzpicture}[scale=0.8]
  % Osy
  \draw[->, thick] (-1,0)--(5,0) node[right]{$x$};
  \draw[->, thick] (0,-1)--(0,4) node[above]{$y$};
  
  % Přímky
  \draw[thick, pastelorange!90!black] (-0.5,2.25)--(4.5,0.75) node[right, black]{$a_1x + b_1y = c_1$};
  \draw[thick, pastelgreen!70!black] (-0.5,0.5)--(4.5,3.5) node[right, black]{$a_2x + b_2y = c_2$};
  
  % Průsečík
  \filldraw[orange] (1.45,1.66) circle (2pt) node[below right] {$[x_0;y_0]$};
\end{tikzpicture}
\end{center}


Podle vzájemné polohy těchto přímek mohou nastat tři základní situace:

\begin{enumerate}[label=\arabic*.]
  \item \textbf{Průsečík existuje právě jeden} — přímky se protínají v jednom bodě.  
  Soustava má \important{jediné řešení}.

  \item \textbf{Přímky se kryjí} — každá hodnota, která vyhovuje jedné rovnici, vyhovuje i druhé.  
  Soustava má \important{nekonečně mnoho řešení}.

  \item \textbf{Přímky jsou rovnoběžné} — žádný bod nevyhovuje oběma rovnicím současně.  
  Soustava \important{nemá řešení}.
\end{enumerate}

\begin{center}
\begin{tikzpicture}[scale=0.75]
  % 1. Jedno řešení
  \begin{scope}
    \draw[->] (-1,0)--(3,0);
    \draw[->] (0,-1)--(0,3);
    \draw[pastelorange!90!black, thick] (-0.5,2)--(2.5,0.5);
    \draw[pastelgreen!70!black, thick] (-0.5,0.5)--(2.5,2);
    \filldraw[orange] (1,1.25) circle (2pt);
    \node at (1,-1.3) {\small jedno řešení};
  \end{scope}

  % 2. Nekonečně mnoho řešení
  \begin{scope}[xshift=5cm]
    \draw[->] (-1,0)--(3,0);
    \draw[->] (0,-1)--(0,3);
    \draw[pastelorange!90!black, thick] (-0.5,0.5)--(2.5,2);
    \draw[pastelgreen!70!black, thick, dashed] (-0.5,0.5)--(2.5,2);
    \node at (1,-1.3) {\small nekonečně mnoho řešení};
  \end{scope}

  % 3. Žádné řešení
  \begin{scope}[xshift=10cm]
    \draw[->] (-1,0)--(3,0);
    \draw[->] (0,-1)--(0,3);
    \draw[pastelorange!90!black, thick] (-0.5,0.5)--(2.5,2);
    \draw[pastelgreen!70!black, thick] (-0.5,0)--(2.5,1.5);
    \node at (1,-1.3) {\small žádné řešení};
  \end{scope}
\end{tikzpicture}
\end{center}

V dalších kapitolách si ukážeme tři základní metody řešení soustav lineárních rovnic:
\begin{enumerate}[label=\arabic*.]
    \item \textbf{Sčítací metodou} – úpravou a sečtením rovnic odstraníme jednu neznámou.
    \item \textbf{Dosazovací metodou} – z jedné rovnice vyjádříme neznámou a dosadíme do druhé.
    \item \textbf{Substituční metodou} – zavedením nové proměnné zjednodušíme zápis soustavy;  
    u lineárních rovnic se však s touto metodou setkáme jen výjimečně – častěji ji využijeme  
    při řešení \important{kvadratických} nebo \important{exponenciálních} rovnic, kde pomáhá odstranit složitější členy.
    \item \textbf{Grafickou metodou} – nalezením průsečíku přímek v souřadnicové rovině.
\end{enumerate}




\subsection{Metody řešení soustav}

\subsubsection*{Sčítací metoda}

Tato metoda se používá zejména tehdy, když všechny koeficienty \(a_1, a_2, b_1, b_2\) jsou nenulové.  
Princip spočívá v tom, že soustavu upravíme tak, aby se jedna z neznámých po sečtení (nebo odečtení) rovnic \textbf{vyrušila}.  
Tím získáme rovnici pouze s jednou neznámou, kterou snadno určíme.

\begin{align*}
x - 5y &= 12,\\
x - 4y &= 15.
\end{align*}

První rovnici vynásobíme \(-1\), abychom získali opačný koeficient u neznámé \(x\):

\begin{align*}
-x + 5y &= -12,\\ 
x - 4y &= 15.
\end{align*}

Když již máme tuto podobu soustavy, můžeme rovnice sečíst:

\[
(-x + 5y) + (x - 4y) = -12 + 15.
\]

Tím dostaneme:

\[
-x + x + 5y - 4y = 3,
\]

a po úpravě:

\[
y = 3.
\]

Tuto hodnotu nyní dosadíme do jedné z původních rovnic, například do první:

\begin{align*}
x - 5y &= 12,\\
x - 5\cdot 3 &= 12,\\
x - 15 &= 12 \quad | +15,\\
x &= 27.
\end{align*}

\textbf{Řešení soustavy:} \quad \(x = 27,\; y = 3\)

\medskip
\textbf{Množina řešení:} \quad \(K = \{[27;3]\}\)

\medskip
Vidíme tedy, že po vhodné úpravě jedné rovnice a následném sečtení obou jsme získali přímou hodnotu jedné z neznámých.  
Tento postup je velmi efektivní zejména tehdy, když se koeficienty jednotlivých neznámých dají snadno převést na opačné hodnoty.

\begin{examplebox}
\textbf{Sčítací metoda}

Řešme soustavu rovnic:
\[
\begin{cases}
3x + 2y = -4,\\
2x - 3y = 19.
\end{cases}
\]

\textit{Postup řešení:}
\begin{align*}
3x + 2y &= -4 \quad | \cdot 3,\\
2x - 3y &= 19 \quad | \cdot 2,\\[-4pt]
\multicolumn{2}{c}{\rule{4cm}{0.3pt}} \\[-2pt]
9x + 6y &= -12,\\
4x - 6y &= 38,\\[-4pt]
\multicolumn{2}{c}{\rule{4cm}{0.3pt}} \\[-2pt]
13x &= 26 \quad | :13,\\
x &= 2.
\end{align*}

Dosadíme do první rovnice:
\begin{align*}
3x + 2y &= -4,\\
3\cdot2 + 2y &= -4,\\
6 + 2y &= -4 \quad | -6,\\
2y &= -10 \quad | :2,\\
y &= -5.
\end{align*}

\textbf{Řešení soustavy:} \quad \(x = 2,\; y = -5\)

\medskip
\textbf{Množina řešení:} \quad \(K = \{[2;-5]\}\)
\end{examplebox}



\subsubsection*{Dosazovací metoda}

Tato metoda je založena na tom, že z jedné rovnice vyjádříme jednu neznámou pomocí druhé  
a tento výraz \textbf{dosadíme do druhé rovnice}. Tím získáme opět jedinou rovnici s jednou neznámou.

\begin{align*}
x - 5y &= 12,\\
x - 4y &= 15.
\end{align*}

Z první rovnice vyjádříme neznámou \(x\):

\[
x = 12 + 5y.
\]

Tento výraz dosadíme do druhé rovnice místo \(x\):

\begin{align*}
(12 + 5y) - 4y &= 15,\\
12 + y &= 15.
\end{align*}

Po úpravě dostaneme:

\[
y = 3.
\]

Nyní dosadíme tuto hodnotu do výrazu \(x = 12 + 5y\):

\begin{align*}
x &= 12 + 5\cdot 3\\
x &= 12 + 15\\
x &= 27
\end{align*}

\textbf{Řešení soustavy:} \quad \(x = 27,\; y = 3\)

\medskip
\textbf{Množina řešení:} \quad \(K = \{[27;3]\}\)

\medskip
Vidíme tedy, že po vyjádření jedné neznámé z první rovnice a jejím dosazení do druhé  
jsme získali přímé vyjádření druhé neznámé. Tento postup bývá obzvláště vhodný, pokud má některá z rovnic jednoduchý koeficient u jedné z neznámých (např. \(1\) nebo \(-1\)).

\begin{examplebox}
\textbf{Dosazovací metoda:}

Řešme následující soustavu rovnic:
\[
\begin{cases}
2x + 7y = 0,\\
5y = -3.
\end{cases}
\]

\textit{Postup řešení:}
\begin{align*}
  2x + 7y &= 0\\
  5y &= -3 \quad \Rightarrow y=-\frac{3}{5}\\[-4pt]
  \multicolumn{2}{c}{\rule{4cm}{0.3pt}} \\[-2pt]
  2x + 7y &= 0\\
  2x + 7\left(-\frac{3}{5}\right) &= 0\\
  2x - \frac{21}{5} &= 0 \quad | +\frac{21}{5}\\
  2x &= \frac{21}{5} \quad | :2\\
  x &= \frac{21}{10}
\end{align*}

\textbf{Řešení soustavy:} \quad \(x = \tfrac{21}{10},\; y = -\tfrac{3}{5}\)

\medskip
\textbf{Množina řešení:} \quad \(K = \left\{\left[\tfrac{21}{10}; -\tfrac{3}{5}\right]\right\}\)
\end{examplebox}

\newpage
\begin{examplebox}
\textbf{Nekonečně mnoho řešení (rovnice jsou násobky)}

Řešme soustavu:
\[
\begin{cases}
x + 2y = 3,\\
2x + 4y = 6.
\end{cases}
\]

\textit{Postup: vynásobíme první rovnici $2$ a odečteme od druhé.}
\begin{align*}
x + 2y &= 3 \quad |\cdot 2,\\
2x + 4y &= 6,\\[-4pt]
\multicolumn{2}{c}{\rule{4cm}{0.3pt}}\\[-2pt]
2x + 4y &= 6,\\
2x + 4y &= 6 \quad | \cdot(-1),\\[-4pt]
\multicolumn{2}{c}{\rule{4cm}{0.3pt}}\\[-2pt]
2x + 4y &= 6,\\
-2x - 4y &= -6 \\[-4pt]
\multicolumn{2}{c}{\rule{4cm}{0.3pt}}\\[-2pt]
0 &= 0.
\end{align*}

Dostali jsme tautologii \(0=0\) — druhá rovnice je \textbf{lineární kombinací} první, takže soustava má \textbf{nekonečně mnoho řešení}: všechny body na přímce \(x + 2y = 3\).

\medskip
\textbf{Zápis množiny řešení:}
\newline
Klasický zápis:
\[
K=\emptyset
\]
nebo parametricky
\[
K = \bigl\{\, [\,3-2t;\; t\,] \;\big|\; t\in\mathbb{R} \,\bigr\}.
\]
\end{examplebox}



\newpage
\section{Soustavy lineárních nerovnic}

Zatímco \textbf{soustavy rovnic} hledají takové hodnoty, které \emph{vyhovují všem rovnicím zároveň} a tedy představují \textbf{průsečík přímek},  
\textbf{soustavy nerovnic} určují \textbf{množinu všech bodů}, které \emph{splňují všechny nerovnosti současně}.  

V rovině proto místo jednoho bodu často získáváme \textbf{oblast} – například trojúhelník, pás nebo polorovinu.

\begin{definitionbox}
\textbf{Soustava lineárních nerovnic o dvou neznámých} je dvojice (nebo více) nerovnic tvaru
\[
\begin{cases}
a_1x + b_1y \,\square\, c_1,\\[4pt]
a_2x + b_2y \,\square\, c_2,\\[2pt]
\end{cases}
\]
kde \(a_1, b_1, c_1, a_2, b_2, c_2  \in \mathbb{R}\) a symbol \(\square \in \{<,\,>,\,\le,\,\ge\}\).  

\textbf{Řešením soustavy} rozumíme všechny dvojice \([x_0;y_0]\), které splňují \emph{všechny} nerovnosti současně.  
Množinu všech řešení značíme obvykle písmenem \(K\).
\end{definitionbox}

\noindent
Z geometrického hlediska představuje každá z těchto nerovnic jednu \textbf{polorovinu},  
a řešení soustavy je jejich \textbf{průnik}.

\begin{notebox}
Každá lineární nerovnice o dvou neznámých představuje v rovině \textbf{polorovinu},  
jejíž hranici tvoří přímka \(a_i x + b_i y = c_i\).  
Pokud je nerovnost ostrá (\(<\) nebo \(>\)), hranice do řešení \textbf{nepatří};  
pokud je neostrá (\(\le\) nebo \(\ge\)), hranice je \textbf{součástí} řešení.
\end{notebox}

\begin{center}
\begin{tikzpicture}[scale=0.8]
  % Osy
  \draw[->, thick] (-1,0)--(5,0) node[right]{$x$};
  \draw[->, thick] (0,-1)--(0,4) node[above]{$y$};
  % Přímka
  \draw[thick, pastelblue!80!black] (-0.5,1.5)--(4.5,3) node[right, black]{$x + 2y = 4$};
  % Vybarvená oblast
  \fill[pastelblue!30, opacity=0.4] (-1,-1)--(5,-1)--(5,3.15)--(-1,1.35)--cycle;
  \node[below right] at (1,0.5) {\small $x + 2y < 4$};
\end{tikzpicture}
\end{center}

\noindent
Vidíme, že každá lineární nerovnice rozdělí rovinu na dvě části – jednu, která danou nerovnici splňuje, a druhou, která ji nesplňuje.  
Při řešení soustavy hledáme oblast, která vyhovuje všem nerovnicím současně.

\begin{examplebox}
\textbf{Příklad:}

Řešme soustavu nerovnic:
\[
\begin{cases}
x + y \le 4,\\
x - y \ge 0.
\end{cases}
\]

\textbf{Řešení:}

1. Přepíšeme nerovnice do tvaru funkcí:
\[
y \le -x + 4, \qquad y \le x.
\]
2. Každá z těchto nerovnic představuje oblast pod (nebo nad) danou přímkou.

\begin{center}
\begin{tikzpicture}[scale=0.9]
  % Osy
  \draw[->, thick] (-1,0)--(5,0) node[right]{$x$};
  \draw[->, thick] (0,-1)--(0,5) node[above]{$y$};
  
  % Přímky
  \draw[thick, pastelblue!80!black] (-0.5,4.5)--(4.5,-0.5) node[right]{$y=-x+4$};
  \draw[thick, pastelgreen!70!black] (-0.5,-0.5)--(4.5,4.5) node[right]{$y=x$};

  % Vybarvení oblasti
  \fill[pastelorange!30, opacity=0.4] (0,0)--(2,2)--(4,0)--cycle;
\end{tikzpicture}
\end{center}

\textbf{Výsledek:}  
Řešením soustavy je trojúhelníková oblast ohraničená přímkami \(y = -x + 4\), \(y = x\) a osou \(x\).
\end{examplebox}

\begin{notebox}
\textbf{Shrnutí:}  
Při řešení soustav nerovnic hledáme oblast, která splňuje všechny podmínky zároveň.  
Typicky jde o průnik několika polorovin – oblast může mít tvar trojúhelníku, čtyřúhelníku nebo pruhu mezi dvěma rovnoběžkami.
\end{notebox}


\newpage
\section{Jedna lineární rovnice o dvou neznámých}

Při řešení lineárních rovnic jsme zatím hledali \textbf{jedinou hodnotu neznámé}, která rovnici splňuje.  
Nyní se zaměříme na rovnice, které obsahují \textbf{dvě neznámé} – například \(x\) a \(y\).  
Tyto rovnice již nemají jediné řešení, ale \textbf{nekonečně mnoho řešení} tvořících určitý geometrický útvar v rovině.

\begin{definitionbox}
\textbf{Lineární rovnice o dvou neznámých} je každá rovnice, kterou lze převést do tvaru
\[
a x + b y = c,
\]
kde \(a, b, c \in \mathbb{R}\) a alespoň jeden z koeficientů \(a\) nebo \(b\) je nenulový.
\end{definitionbox}

Každé řešení této rovnice představuje \textbf{uspořádanou dvojici čísel} \([x_0; y_0]\), která po dosazení splňuje vztah \(a x_0 + b y_0 = c\).  
Množinu všech těchto dvojic označujeme \(K\).

\medskip
Zatímco lineární rovnice o jedné neznámé má řešení v podobě \textbf{bodů na číselné ose},  
rovnice o dvou neznámých má řešení v podobě \textbf{bodů v rovině}.  
Tato množina bodů tvoří \textbf{přímku}.

\begin{notebox}
Každá lineární rovnice o dvou neznámých představuje v rovině \textbf{přímku}.  
Každý bod této přímky odpovídá jednomu řešení rovnice.
\end{notebox}

\medskip
Rovnici \(a x + b y = c\) můžeme vyjádřit také ve tvaru
\[
y = -\frac{a}{b}x + \frac{c}{b},
\]
pokud \(b \neq 0\). Tento tvar známe z lineární funkce \(y = kx + q\),  
kde směrnice přímky je \(k = -\frac{a}{b}\) a průsečík s osou \(y\) je \(q = \frac{c}{b}\).

\begin{examplebox}
\textbf{Příklad:}

Uvažujme rovnici
\[
2x + 3y = 6.
\]

Vyjádříme \(y\) pomocí \(x\):
\[
3y = -2x + 6 \quad \Rightarrow \quad y = -\frac{2}{3}x + 2.
\]

Rovnice tedy představuje přímku se směrnicí \(k = -\frac{2}{3}\)  
a průsečíkem s osou \(y\) v bodě \([0;2]\).

\begin{center}
\begin{axisgraph}[
    xmin=-1, xmax=6,
    ymin=-1, ymax=4,
    axis lines=middle,
    axis line style={thick},
    xlabel={$x$}, ylabel={$y$},
    width=0.75\linewidth,
    height=6cm,
    samples=100,
    domain=-1:6,
    grid=none,
    title={Graf rovnice $2x + 3y = 6$}
]
    \addplot[thick, pastelblue!80!black, domain=-1:6]{-2/3*x + 2};
    \addplot[only marks, mark=*, pastelblue!80!black] coordinates {(0,2) (3,0)};
    \node[above right, pastelblue!80!black] at (axis cs:3,0) {$[3;0]$};
    \node[right, pastelblue!80!black] at (axis cs:0,2) {$[0;2]$};
\end{axisgraph}
\end{center}

\textbf{Výsledek:}  
Množinu všech řešení tvoří přímka procházející body \([0;2]\) a \([3;0]\).

\[
K = \{[x;y] \in \mathbb{R}^2 ; 2x + 3y = 6\}.
\]

Ekvivalentně můžeme množinu řešení vyjádřit i tak, že jedna proměnná je volná a druhou dopočítáme z rovnice.  
Z výrazu \(y = -\frac{2}{3}x + 2\) tedy platí:
\[
K = \left\{[x;\,-\tfrac{2}{3}x + 2] ; x \in \mathbb{R}\right\}.
\]

\medskip
Tento zápis zdůrazňuje, že každé reálné číslo \(x\) určuje právě jeden bod \([x;y]\) na přímce —  
množinu všech takových bodů pak tvoří celá přímka, která je řešením dané rovnice.

\end{examplebox}

\medskip
Z geometrického hlediska tedy platí, že každá rovnice o dvou neznámých popisuje všechny body, které leží na určité přímce v rovině.  
Pokud k této rovnici přidáme druhou (tedy vytvoříme soustavu), hledáme \textbf{průsečík} těchto přímek – právě ten odpovídá společnému řešení.

\begin{notebox}
\textbf{Shrnutí:}
\begin{itemize}
    \item Rovnice o jedné neznámé → řešením jsou body na číselné ose.
    \item Rovnice o dvou neznámých → řešením je přímka v rovině.
    \item Soustava dvou rovnic → řešením je průsečík dvou přímek.
\end{itemize}
\end{notebox}





\newpage
\section{Lineární rovnice s absolutní hodnotou}

V předchozích kapitolách jsme řešili rovnice, ve kterých se neznámá vyskytovala v lineárním tvaru, tedy bez mocnin či jiných operací.  
Nyní se setkáme s případy, kdy je neznámá součástí \textbf{absolutní hodnoty}.  
Takové rovnice se sice tváří jednoduše, ale vyžadují určitou pozornost, protože \textbf{absolutní hodnota může skrývat dvě různé situace}.  

\medskip
Absolutní hodnota čísla vyjadřuje jeho \textbf{vzdálenost od nuly} na číselné ose —  
nezáleží tedy na tom, zda je číslo kladné, nebo záporné, výsledná hodnota je vždy nezáporná.  
Například \(|3| = 3\) a \(|-3| = 3\).

\begin{definitionbox}
\textbf{Absolutní hodnota reálného čísla} \(a\) je definována takto:
\[
|a| =
\begin{cases}
a, & \text{je-li } a \ge 0,\\[4pt]
-a, & \text{je-li } a < 0.
\end{cases}
\]
\end{definitionbox}

\noindent
Z této definice můžeme odvodit několik základních vlastností,  
které nám usnadní práci při výpočtech s absolutní hodnotou.

\begin{notebox}
\textbf{Vlastnosti absolutní hodnoty:}
\[
\begin{aligned}
|a| &\ge 0, \\[4pt]
|a| &= |-a|, \\[4pt]
|a \cdot b| &= |a| \cdot |b|, \\[4pt]
\left|\frac{a}{b}\right| &= \frac{|a|}{|b|}, \quad \text{je-li } b \neq 0.
\end{aligned}
\]
\end{notebox}

\noindent
Z geometrického hlediska tedy absolutní hodnota představuje \textbf{vzdálenost bodu} s souřadnicí \(a\) od počátku souřadnicového systému.  
Proto je výsledkem vždy nezáporné číslo.

\begin{center}
\begin{tikzpicture}[scale=1]
  \draw[->, thick] (-4,0)--(5,0) node[right]{$x$};
  \foreach \x in {-3,-2,-1,0,1,2,3,4}
    \draw (\x,0.1)--(\x,-0.1) node[below]{\small \x};
  \draw[thick, pastelblue, fill=pastelblue!30] (3,0) circle (2pt);
  \draw[thick, pastelblue, fill=pastelblue!30] (-3,0) circle (2pt);
  \draw[<->, gray!70] (-3,0.4)--(0,0.4) node[midway, above]{$3$};
  \draw[<->, gray!70] (0,0.4)--(3,0.4) node[midway, above]{$3$};
\end{tikzpicture}
\end{center}

\noindent
Například čísla \(-3\) a \(3\) mají od nuly stejnou vzdálenost — právě tuto vzdálenost zachycuje absolutní hodnota.  

\begin{notebox}
Absolutní hodnota čísla představuje jeho vzdálenost od nuly,  
a proto je výsledkem vždy číslo \textbf{nezáporné}.
\end{notebox}

\medskip
Pokud se absolutní hodnota objeví v rovnici, musíme rozlišit,  
\textbf{zda je výraz uvnitř kladný, nebo záporný}.  
To znamená, že každá taková rovnice vede na \textbf{dvě lineární rovnice},  
které řešíme samostatně.

\begin{definitionbox}
\textbf{Obecný tvar lineární rovnice s absolutní hodnotou:}
\[
|ax + b| = c,
\]
kde \(a, b, c \in \mathbb{R}\).

\textbf{Řešení:}
\begin{itemize}
  \item pokud \(c < 0\), rovnice nemá řešení (absolutní hodnota nemůže být záporná);
  \item pokud \(c \ge 0\), řešíme dvojici rovnic:
  \[
  \begin{cases}
  ax + b = c,\\[4pt]
  ax + b = -c.
  \end{cases}
  \]
\end{itemize}
\end{definitionbox}

\noindent
Tyto dvě rovnice řešíme obvyklým způsobem a jejich řešení pak spojíme do jedné množiny.

\medskip
Nyní si ukážeme základní příklad, na kterém si princip řešení vysvětlíme krok za krokem.

\begin{examplebox}
\textbf{Příklad:}  
Vyřešte rovnici
\[
|x - 5| = 2.
\]

\textbf{1. Algebraický postup:}

Podle definice absolutní hodnoty rozlišíme dva případy:
\[
\begin{cases}
x - 5 = 2,\\[4pt]
x - 5 = -2.
\end{cases}
\]

\textbf{Výsledek:}

\[
K = \{3, 7\},
\begin{cases}
x = 7,\\[4pt]
x = 3.
\end{cases}
\]

\medskip
\textbf{2. Geometrická představa:}

Absolutní hodnota vyjadřuje vzdálenost.  
Rovnice \(|x - 5| = 2\) říká, že hledáme všechna čísla, jejichž vzdálenost od čísla \(5\) je rovna \(2\).

\begin{center}
\begin{tikzpicture}[scale=0.9]
  \draw[->, thick] (-1,0)--(9,0) node[right]{$x$};
  \foreach \x in {0,1,...,8}
    \draw (\x,0.1)--(\x,-0.1) node[below]{\small \x};
  % body
  \filldraw[pastelblue] (5,0) circle (2pt) node[above=3pt]{$5$};
  \filldraw[pastelorange] (3,0) circle (2pt) node[above=3pt]{$3$};
  \filldraw[pastelorange] (7,0) circle (2pt) node[above=3pt]{$7$};
  % vzdálenosti
  \draw[<->, gray!70] (3,0.6)--(5,0.6) node[midway, above]{$2$};
  \draw[<->, gray!70] (5,0.6)--(7,0.6) node[midway, above]{$2$};
\end{tikzpicture}
\end{center}

Z grafu je zřejmé, že od bodu \(5\) leží dva body ve vzdálenosti \(2\) – tedy \(x = 3\) a \(x = 7\).  
Tyto hodnoty odpovídají i algebraickému řešení.

\end{examplebox}

\noindent
Vidíme tedy, že algebraický a geometrický přístup dávají stejný výsledek –  
zatímco algebraické řešení vychází z definice absolutní hodnoty,  
geometrická představa zdůrazňuje význam jako \textbf{vzdálenosti na číselné ose}.

\begin{notebox}
\textbf{Shrnutí:}
\begin{itemize}
  \item Absolutní hodnota představuje vzdálenost od nuly, a proto je vždy nezáporná.
  \item Rovnici \(|ax + b| = c\) řešíme rozborem dvou případů:
  \[
  \begin{cases}
  ax + b = c,\\[4pt]
  ax + b = -c.
  \end{cases}
  \]
  \item Geometricky hledáme všechny body, jejichž vzdálenost od bodu s souřadnicí \(-\tfrac{b}{a}\) je rovna \(\tfrac{c}{|a|}\).
\end{itemize}
\end{notebox}


Mějme například tuto rovnici:
\[
|x - 2| + |2x - 8| = 5.
\]

Tento příklad je složitější, protože obsahuje \textbf{dvě absolutní hodnoty}.  
V takovém případě už nestačí jen rozdělit na dva případy –  
musíme nejprve zjistit, \textbf{v jakých intervalech mají výrazy uvnitř absolutních hodnot kladné nebo záporné hodnoty}.  

\medskip
Každá absolutní hodnota se mění v bodě, kde je její vnitřní výraz nulový:

\[
\begin{aligned}
x - 2 &= 0 \quad \Rightarrow \quad x = 2,\\
2x - 8 &= 0 \quad \Rightarrow \quad x = 4.
\end{aligned}
\]

Těmito body si tedy rozdělíme číselnou osu na tři intervaly:
\[
(-\infty, 2), \quad \langle 2, 4), \quad \langle 4, \infty).
\]

V každém z nich má výraz uvnitř absolutní hodnoty jiná znaménka,  
a proto musíme v každém intervalu rovnici rozepsat zvlášť.

\begin{center}
\begin{tikzpicture}[scale=1]
  \draw[->, thick] (-3,0)--(8,0) node[right]{$x$};
  \foreach \x in {-2,-1,0,1,2,3,4,5,6,7}
    \draw (\x,0.1)--(\x,-0.1) node[below]{\small \x};
  % body
  \filldraw[pastelorange!80!black] (2,0) circle (2pt) node[above=3pt]{$2$};
  \filldraw[pastelorange!80!black] (4,0) circle (2pt) node[above=3pt]{$4$};
  % intervals
  \node[above, pastelorange!80!black] at (0.2,0.3) {\small $x<2$};
  \node[above, pastelorange!80!black] at (3,0.3) {\small $2\le x<4$};
  \node[above, pastelorange!80!black] at (6,0.3) {\small $x\ge4$};
\end{tikzpicture}
\end{center}

\textbf{1. interval:} \(x < 2\)

Oba výrazy jsou záporné, proto oba měníme na opačné:
\[
-(x - 2) - (2x - 8) = 5.
\]
\[
-x + 2 - 2x + 8 = 5 \quad \Rightarrow \quad -3x + 10 = 5 \quad \Rightarrow \quad x = \tfrac{5}{3}.
\]
Tento výsledek patří do intervalu \((- \infty, 2)\), takže jej přijímáme.

\[
\Rightarrow \; x_1 = \tfrac{5}{3}.
\]

\medskip
\textbf{2. interval:} \(2 \le x < 4\)

Zde první výraz \((x - 2)\) je kladný, druhý \((2x - 8)\) je záporný.
\[
(x - 2) - (2x - 8) = 5.
\]
\[
x - 2 - 2x + 8 = 5 \quad \Rightarrow \quad -x + 6 = 5 \quad \Rightarrow \quad x = 1.
\]
Tento výsledek však do intervalu \(\langle 2, 4)\) nepatří, proto ho \textbf{nepřijímáme.}

\medskip
\textbf{3. interval:} \(x \ge 4\)

Oba výrazy jsou kladné:
\[
(x - 2) + (2x - 8) = 5.
\]
\[
3x - 10 = 5 \quad \Rightarrow \quad x = 5.
\]
Tento výsledek do intervalu \(\langle 4, \infty)\) patří, takže ho přijímáme.

\[
\Rightarrow \; x_2 = 5.
\]

\medskip
\textbf{Výsledek:}
\[
K = \left\{\tfrac{5}{3},\,5\right\}.
\]


\noindent
Vidíme, že u rovnic s více absolutními hodnotami musíme nejprve určit intervaly,  
ve kterých se výrazy uvnitř absolutních hodnot nemění ve znaménku.  
Teprve poté lze správně sestavit jednotlivé lineární rovnice a určit všechna řešení.

\begin{notebox}
\textbf{Postup řešení rovnic s více absolutními hodnotami:}
\begin{enumerate}[label=\arabic*.]
  \item Najdi body, ve kterých se vnitřní výrazy rovnají nule.
  \item Rozděl číselnou osu podle těchto bodů na intervaly.
  \item V každém intervalu nahraď absolutní hodnoty příslušnými výrazy se správným znaménkem.
  \item Řeš vzniklé lineární rovnice a ověř, zda výsledky do daných intervalů patří.
\end{enumerate}
\end{notebox}

\begin{examplebox}
\textbf{Příklad:}  
Vyřešte rovnici
\[
|x - 2| = |x + 1|.
\]

Nulové body jednotlivých absolutních hodnot jsou \(x = 2\) a \(x = -1\).  
Číselnou osu si tedy rozdělíme na tři intervaly:
\[
(-\infty, -1), \quad \langle -1, 2), \quad \langle 2, \infty).
\]

\newpage
\textbf{1. interval:} \(x < -1\)

Oba výrazy jsou záporné:
\[
-(x - 2) = -(x + 1) \quad \Rightarrow \quad -x + 2 = -x - 1.
\]
Tato rovnost neplatí (žádné řešení).

\medskip
\textbf{2. interval:} \(-1 \le x < 2\)

Zde první výraz \((x - 2)\) je záporný a druhý \((x + 1)\) kladný:
\[
-(x - 2) = x + 1 \quad \Rightarrow \quad -x + 2 = x + 1 \quad \Rightarrow \quad 2x = 1 \quad \Rightarrow \quad x = \tfrac{1}{2}.
\]
Tento výsledek do intervalu patří, tedy jej přijímáme.

\medskip
\textbf{3. interval:} \(x \ge 2\)

Oba výrazy jsou kladné:
\[
x - 2 = x + 1 \quad \Rightarrow \quad -2 = 1,
\]
což není pravda (žádné řešení).

\medskip
\textbf{Výsledek:}
\[
K = \left\{\tfrac{1}{2}\right\}.
\]
\end{examplebox}




\subsection{Lineární nerovnice s absolutní hodnotou}

Stejně jako u rovnic, i u nerovnic se může absolutní hodnota objevit v různých tvarech.  
Připomeňme si, že \textbf{absolutní hodnota vyjadřuje vzdálenost od nuly} –  
a proto je vždy \textbf{nezáporná}.  

\begin{definitionbox}
\textbf{Absolutní hodnota reálného čísla} \(a\) je definována:
\[
|a| =
\begin{cases}
a, & a \ge 0,\\[4pt]
-a, & a < 0.
\end{cases}
\]

Z toho plyne, že:
\[
|a| \ge 0 \quad \text{pro všechna } a \in \mathbb{R}.
\]
\end{definitionbox}

V nerovnicích absolutní hodnota určuje, jak \textbf{daleko} od určitého bodu může být hodnota proměnné.  
To znamená, že místo jednoho nebo dvou konkrétních řešení hledáme celou \textbf{množinu čísel}, která vyhovují dané podmínce.


\begin{notebox}
\textbf{Základní princip:}

Mějme nerovnici tvaru
\[
|x - a| < b,
\]
kde \(b > 0\).

\medskip
Tato nerovnice říká, že „vzdálenost mezi číslem \(x\) a číslem \(a\) je menší než \(b\)“.  
Tedy \(x\) leží v intervalu:
\[
a - b < x < a + b.
\]

Naopak, pro nerovnici
\[
|x - a| > b
\]
platí, že \(x\) leží \textbf{mimo tento interval}, tedy:
\[
x < a - b \quad \text{nebo} \quad x > a + b.
\]
\end{notebox}


\begin{examplebox}
\textbf{Příklad 1:}  
Vyřešte nerovnici
\[
|x - 3| < 2.
\]

\textbf{Řešení:}  
Tato nerovnice říká, že \(x\) má být od čísla \(3\) vzdálen méně než \(2\).  
Podle definice tedy:
\[
-2 < x - 3 < 2.
\]

Po přičtení čísla \(3\):
\[
1 < x < 5.
\]

\textbf{Výsledek:}
\[
K = \langle 1; 5 \rangle.
\]

\begin{center}
\begin{tikzpicture}[scale=1]
  \draw[->, thick] (-1,0)--(7,0) node[right]{$x$};
  \foreach \x in {0,1,2,3,4,5,6}
    \draw (\x,0.1)--(\x,-0.1) node[below]{\small \x};
  % interval
  \draw[thick, pastelorange, line width=1pt] (1,0)--(5,0);
  % otevřené body
  \draw[thick, pastelorange!80!black] (1,0) circle (2pt);
  \draw[thick, pastelorange!80!black] (5,0) circle (2pt);
  \node[above, pastelorange!80!black] at (3,0.3) {\small $|x-3|<2$};
\end{tikzpicture}
\end{center}
\end{examplebox}


\begin{examplebox}
\textbf{Příklad 2:}  
Vyřešte nerovnici
\[
|x - 1| \ge 3.
\]

\textbf{Řešení:}  
Tato nerovnice znamená, že \(x\) je od čísla \(1\) vzdálen alespoň o \(3\).  
Podle definice tedy:
\[
x - 1 \ge 3 \quad \text{nebo} \quad x - 1 \le -3.
\]

Z toho:
\[
x \ge 4 \quad \text{nebo} \quad x \le -2.
\]

\textbf{Výsledek:}
\[
K = (-\infty, -2\rangle \cup \langle 4, \infty).
\]

\begin{center}
\begin{tikzpicture}[scale=1]
  \draw[->, thick] (-5,0)--(7,0) node[right]{$x$};
  \foreach \x in {-4,-3,-2,-1,0,1,2,3,4,5,6}
    \draw (\x,0.1)--(\x,-0.1) node[below]{\small \x};
  % intervals
  \draw[thick, pastelblue, line width=1pt, ->] (4,0)--(6.5,0);
  \draw[thick, pastelblue, line width=1pt, <-] (-4.5,0)--(-2,0);
  % body
  \filldraw[pastelblue!80!black] (4,0) circle (2pt);
  \filldraw[pastelblue!80!black] (-2,0) circle (2pt);
  \node[above, pastelblue!80!black] at (1,0.3) {\small $|x-1|\ge3$};
\end{tikzpicture}
\end{center}
\end{examplebox}

\newpage
\begin{notebox}
\textbf{Shrnutí:}

\[
\begin{array}{rcll}
|x - a| < b & \Rightarrow & a - b < x < a + b & \text{(vzdálenost menší než $b$)}\\[4pt]
|x - a| \le b & \Rightarrow & a - b \le x \le a + b &\\[4pt]
|x - a| > b & \Rightarrow & x < a - b \;\text{nebo}\; x > a + b &\\[4pt]
|x - a| \ge b & \Rightarrow & x \le a - b \;\text{nebo}\; x \ge a + b &
\end{array}
\]


\medskip
\noindent
Tyto vzorce platí pro všechna \(b > 0\).  
Pokud by \(b < 0\), nerovnice nemá žádné řešení, protože absolutní hodnota je vždy nezáporná.
\end{notebox}



\newpage
\subsection{Lineární rovnice s parametrem}

V některých rovnicích se kromě neznámé \(x\) objevuje i další proměnná,  
která není určena — tzv. \textbf{parametr}.  
Parametr označujeme zpravidla písmeny \(m, a, k, p\) apod.  
Mění-li se hodnota parametru, mění se i samotná rovnice,  
a tedy i \textbf{počet nebo tvar řešení}.  

\medskip
Takové rovnice nám umožňují nejen hledat konkrétní řešení,  
ale také zkoumat, \textbf{za jakých podmínek vůbec řešení existuje}.  

\begin{definitionbox}
\textbf{Lineární rovnice s parametrem} je rovnice, ve které se kromě neznámé vyskytuje i další proměnná – parametr.  
Obecný tvar je:
\[
(a + b m)x = c m + d,
\]
kde \(m\) je parametr a \(a, b, c, d \in \mathbb{R}\).
\end{definitionbox}


Rovnici s parametrem řešíme podobně jako běžnou lineární rovnici,  
ale zároveň musíme zohlednit, zda některé hodnoty parametru nevedou k \textbf{nesmyslnému} nebo \textbf{neurčitému} vztahu —  
například k dělení nulou nebo k rovnosti \(0 = 0\).  

\textbf{Rozbor možných situací:}
\begin{itemize}
  \item Pokud se po úpravě objeví tvar \(a x = b\) a \(a \neq 0\),  
        má rovnice \textbf{jedno řešení}.
  \item Pokud \(a = 0\) a zároveň \(b \neq 0\),  
        dostáváme \textbf{nesmyslnou rovnici} (např. \(0 = 5\)) — žádné řešení.
  \item Pokud \(a = 0\) a \(b = 0\),  
        vyjde \textbf{pravdivé tvrzení} \(0 = 0\) — nekonečně mnoho řešení.
\end{itemize}


\textbf{Příklad 1 (komentovaný):}

Uvažujme rovnici
\[
(m - 1)x = 3m + 2.
\]

Vidíme, že parametr \(m\) ovlivňuje koeficient u \(x\).  
Abychom zjistili, kdy rovnice má či nemá řešení,  
rozdělíme úlohu podle hodnoty výrazu \(m - 1\).

\begin{itemize}
  \item \textbf{Případ 1:} \(m - 1 \neq 0\), tj. \(m \neq 1\)

  Potom můžeme rovnici normálně vydělit:
  \[
  x = \frac{3m + 2}{m - 1}.
  \]
  Rovnice má právě \textbf{jedno řešení} pro každé \(m \neq 1\).

  \item \textbf{Případ 2:} \(m - 1 = 0\), tj. \(m = 1\)

  Dosadíme do původní rovnice:
  \[
  (1 - 1)x = 3(1) + 2 \quad \Rightarrow \quad 0 = 5.
  \]
  To je nepravdivé tvrzení – rovnice tedy \textbf{nemá řešení.}
\end{itemize}

\textbf{Shrnutí:}
\[
\begin{array}{l|l}
\text{Hodnota parametru } m & \text{Počet řešení}\\
\hline
m \neq 1 & 1 \text{ řešení},\\
m = 1 & \text{žádné řešení.}
\end{array}
\]

\textbf{Množina řešení:} \quad
\(K_m = \left\{\left[\dfrac{3m+2}{m-1}\right]\right\}, \; m \neq 1.\)


\begin{examplebox}
\textbf{Příklad 2}

Určete, pro jaké hodnoty parametru \(a \in \mathbb{R}\) má rovnice
\[
(a - 2)x = 4a - 8
\]
žádné, jedno nebo nekonečně mnoho řešení.

\textit{Postup:}
\begin{align*}
(a - 2)x &= 4a - 8,\\
\multicolumn{2}{c}{\rule{6cm}{0.3pt}}\\[-2pt]
\text{Pokud } a - 2 \neq 0 &\Rightarrow x = \dfrac{4a - 8}{a - 2},\\
\text{Pokud } a - 2 = 0 &\Rightarrow a = 2,\\
0 \cdot x &= 4\cdot2 - 8 \;\Rightarrow\; 0 = 0.
\end{align*}

Dostáváme tedy:
\[
\begin{array}{l|l}
\text{Hodnota parametru } a & \text{Počet řešení}\\
\hline
a \neq 2 & 1 \text{ řešení},\\
a = 2 & \text{nekonečně mnoho řešení.}
\end{array}
\]

\textbf{Množina řešení:}
\[
K_a =
\begin{cases}
\displaystyle \left\{\left[\frac{4a-8}{a-2}\right]\right\}, & a \neq 2,\\[8pt]
\mathbb{R}, & a = 2.
\end{cases}
\]
\end{examplebox}


Vidíme, že parametrické rovnice nám umožňují nejen vypočítat konkrétní řešení,  
ale i analyzovat chování celé rodiny rovnic v závislosti na hodnotě parametru.  
Takový přístup se později uplatní i při řešení \textbf{soustav lineárních rovnic},  
kde se o počtu řešení rozhoduje podle hodnot determinantů nebo parametrů v matici.





\newpage
\section{Soustava tří LR o třech neznámých}

Zatímco soustava dvou rovnic o dvou neznámých představovala průsečík dvou přímek v rovině,  
soustava tří rovnic o třech neznámých odpovídá hledání \textbf{průsečíku tří rovin v prostoru}.  
Tak jako se dvě přímky mohou protínat, být rovnoběžné nebo splývat,  
mohou se i tři roviny protínat v jednom bodě, v přímce, nebo se vůbec neprotínat.

\begin{definitionbox}
\textbf{Soustava tří lineárních rovnic o třech neznámých} má tvar
\[
\begin{cases}
a_1x + b_1y + c_1z = d_1,\\
a_2x + b_2y + c_2z = d_2,\\
a_3x + b_3y + c_3z = d_3,
\end{cases}
\]
kde \(a_i, b_i, c_i, d_i \in \mathbb{R}\).

\medskip
\textbf{Řešením} soustavy je každá trojice \([x_0; y_0; z_0]\), která vyhovuje všem třem rovnicím zároveň.
\end{definitionbox}

\noindent
Z geometrického hlediska tedy každá rovnice představuje rovinu v prostoru,  
a jejich společný průsečík je řešení celé soustavy.

\begin{center}
\tdplotsetmaincoords{70}{110}

\begin{minipage}{0.47\textwidth}
\begin{tikzpicture}[tdplot_main_coords, scale=0.9, font=\sffamily]

% --- roviny ---
\filldraw[fill=pastelorange!80, draw=pastelorange!80!black, opacity=0.5]
  (-3,0,-3) -- (-3,0,3) -- (3,0,3) -- (3,0,-3) -- cycle; % xy

\filldraw[fill=pastelblue!80, draw=pastelblue!70!black, opacity=0.5]
  (-3,-3,0) -- (-3,3,0) -- (3,3,0) -- (3,-3,0) -- cycle; % xz

\filldraw[fill=pastelgreen!80, draw=pastelgreen!70!black, opacity=0.5]
  (2,-3,-3) -- (2,3,-3) -- (-2,3,3) -- (-2,-3,3) -- cycle; % šikmá

% --- průsečnice ---
\draw[very thick, black] (-3,0,0) -- (3,0,0);
\draw[very thick, black] (0,-3,0) -- (0,3,0);
\draw[very thick, black] (0,-0.75,-4) -- (0,0.75,4);

% --- průsečík ---
\filldraw[red!80!black] (0,0,0) circle (2.5pt);
\node[below right, red!80!black] at (0,0,0) {\small $[x_0;\,y_0;\,z_0]$};

% --- popisek ---
\node[anchor=south west, align=left] (txt) at (4,4,3.5)
  {\small \textbf{Průsečík}\\ tří rovin};
\draw[-latex, thick, red!70!black]
  (txt) to[out=180, in=60] (0.2,0.2,0.1);

\end{tikzpicture}

\vspace{3pt}
\centering
\small\textit{Soustava má jedno řešení – roviny se protínají v jediném bodě.}
\end{minipage}
\hfill
\begin{minipage}{0.47\textwidth}
\begin{tikzpicture}[tdplot_main_coords, scale=0.9, font=\sffamily]

% --- roviny ---
\filldraw[fill=pastelorange!80, draw=pastelorange!80!black, opacity=0.5]
  (-3,-1,-3) -- (-3,1,3) -- (3,1,3) -- (3,-1,-3) -- cycle;

\filldraw[fill=pastelblue!80, draw=pastelblue!70!black, opacity=0.5]
  (-3,-3,0) -- (-3,3,0) -- (3,3,0) -- (3,-3,0) -- cycle;

\filldraw[fill=pastelgreen!80, draw=pastelgreen!70!black, opacity=0.5]
  (3,3,-3) -- (-3,3,-3) -- (-3,-3,3) -- (3,-3,3) -- cycle;

% --- průsečnice (roviny se protínají v přímce) ---
\draw[very thick, black] (-3,0,0) -- (3,0,0);

% --- průsečná přímka (nekonečně mnoho řešení) ---
\draw[ultra thick, red!70!black] (-3,0,0) -- (3,0,0);
\node[below right, red!80!black] at (1.8,0,0) {\small průsečná přímka};

\end{tikzpicture}

\vspace{3pt}
\centering
\small\textit{Soustava má nekonečně mnoho řešení – roviny se protínají v přímce.}
\end{minipage}
\end{center}




\begin{notebox}
\textbf{Z geometrického hlediska:}
\begin{itemize}
  \item Průsečík tří rovin → jedno řešení.
  \item Společná přímka → nekonečně mnoho řešení.
  \item Rovnoběžné nebo nespojitě položené roviny → žádné řešení.
\end{itemize}
\end{notebox}


\newpage
\subsection{Řešení klasickým postupem}

Uvažujme soustavu:
\[
\begin{cases}
x + y + z = 6,\\
2x - y + 3z = 14,\\
-x + 2y + z = 2.
\end{cases}
\]

Z první rovnice vyjádříme \(x\):
\[
x = 6 - y - z.
\]

Dosadíme do zbývajících dvou rovnic:
\[
\begin{cases}
2(6 - y - z) - y + 3z = 14,\\
-(6 - y - z) + 2y + z = 2.
\end{cases}
\]

Po úpravách:
\[
\begin{cases}
12 - 2y - 2z - y + 3z = 14,\\
-6 + y + z + 2y + z = 2.
\end{cases}
\]

\[
\begin{cases}
-3y + z = 2,\\
3y + 2z = 8.
\end{cases}
\]

Tím jsme získali soustavu dvou rovnic o dvou neznámých.  
Sečteme je:
\[
(-3y + z) + (3y + 2z) = 2 + 8 \quad \Rightarrow \quad 3z = 10 \quad \Rightarrow \quad z = \tfrac{10}{3}.
\]

Dosadíme do první rovnice:
\[
-3y + \tfrac{10}{3} = 2 \quad \Rightarrow \quad -3y = -\tfrac{4}{3} \quad \Rightarrow \quad y = \tfrac{4}{9}.
\]

A nakonec do \(x = 6 - y - z\):
\[
x = 6 - \tfrac{4}{9} - \tfrac{10}{3} = \tfrac{40}{9}.
\]

\[
\textbf{Řešení soustavy:} \quad x = \tfrac{40}{9},\; y = \tfrac{4}{9},\; z = \tfrac{10}{3}.
\]

\[
K = \left\{\left[\tfrac{40}{9}; \tfrac{4}{9}; \tfrac{10}{3}\right]\right\}.
\]


\begin{examplebox}
\textbf{Příklad:}  
Vyřešte soustavu tří lineárních rovnic o třech neznámých:
\[
\begin{cases}
2x + y - z = 8,\\
x - 2y + 3z = -4,\\
3x + y + 2z = 10.
\end{cases}
\]

\textbf{Řešení:}

Z první rovnice vyjádříme \(y\):
\[
y = 8 - 2x + z.
\]

Dosadíme do druhé a třetí rovnice:
\[
\begin{cases}
x - 2(8 - 2x + z) + 3z = -4,\\
3x + (8 - 2x + z) + 2z = 10.
\end{cases}
\]

Po úpravě:
\[
\begin{cases}
x - 16 + 4x - 2z + 3z = -4,\\
3x + 8 - 2x + z + 2z = 10.
\end{cases}
\]

\[
\begin{cases}
5x + z = 12,\\
x + 3z = 2.
\end{cases}
\]

Vynásobíme druhou rovnici \(-5\) a přičteme k první:
\[
(5x + z) + (-5x - 15z) = 12 - 10 \quad \Rightarrow \quad -14z = 2 \quad \Rightarrow \quad z = -\tfrac{1}{7}.
\]

Dosadíme \(z = -\tfrac{1}{7}\) do \(x + 3z = 2\):
\[
x + 3\left(-\tfrac{1}{7}\right) = 2 \quad \Rightarrow \quad x = 2 + \tfrac{3}{7} = \tfrac{17}{7}.
\]

Dosadíme \(x, z\) do \(y = 8 - 2x + z\):
\[
y = 8 - 2\cdot\tfrac{17}{7} - \tfrac{1}{7} = \tfrac{56 - 34 - 1}{7} = \tfrac{21}{7} = 3.
\]

\[
\textbf{Řešení:} \quad x = \tfrac{17}{7},\; y = 3,\; z = -\tfrac{1}{7}.
\]

\[
\textbf{Množina řešení:} \quad K = \left\{\left[\tfrac{17}{7};\,3;\,-\tfrac{1}{7}\right]\right\}.
\]
\end{examplebox}



\subsection{Maticový zápis soustavy}

Předchozí postup je funkční, ale zdlouhavý.  
Pro přehlednější zápis a jednodušší práci se zavádí pojem \textbf{matice}.

Každou soustavu lineárních rovnic lze zapsat ve tvaru:
\[
A \cdot X = B,
\]
kde:
\[
A =
\begin{bmatrix}
a_1 & b_1 & c_1\\
a_2 & b_2 & c_2\\
a_3 & b_3 & c_3
\end{bmatrix},
\quad
X =
\begin{bmatrix}
x\\y\\z
\end{bmatrix},
\quad
B =
\begin{bmatrix}
d_1\\d_2\\d_3
\end{bmatrix}.
\]

\noindent
Pro praktické výpočty používáme \textbf{rozšířenou matici}:
\[
\left[
\begin{array}{ccc|c}
a_1 & b_1 & c_1 & d_1\\
a_2 & b_2 & c_2 & d_2\\
a_3 & b_3 & c_3 & d_3
\end{array}
\right].
\]


\subsection{Gaussova eliminační metoda (GEM)}

\textbf{Princip:}  
Pomocí elementárních úprav řádků převedeme rozšířenou matici na jednodušší tvar,  
ze kterého lze snadno dopočítat jednotlivé neznámé.

\begin{notebox}
\textbf{Povolené úpravy:}
\begin{itemize}
  \item násobení řádku nenulovým číslem,
  \item přičtení násobku jednoho řádku k jinému,
  \item prohození dvou řádků.
\end{itemize}
\end{notebox}

\textbf{Cíl:}  
Získat tzv. \textbf{horní trojúhelníkový tvar}, tj. tvar s nulami pod hlavní diagonálou:
\[
\begin{bmatrix}
1 & * & * & *\\
0 & 1 & * & *\\
0 & 0 & 1 & *
\end{bmatrix}.
\]
Z takové matice už snadno určíme \(z\), poté \(y\), nakonec \(x\) — tzv. \textbf{zpětnou substitucí}.

\begin{examplebox}
\textbf{Řešení soustavy tří rovnic:}

Řešme soustavu:
\[
\begin{cases}
2x + y - z = 8,\\
x - 2y + 3z = -4,\\
3x + y + 2z = 10.
\end{cases}
\]

\textit{Z první rovnice vyjádříme } \(y\):
\[
y = 8 - 2x + z.
\]

Dosadíme do zbývajících dvou rovnic:
\begin{align*}
x - 2(8 - 2x + z) + 3z &= -4,\\
3x + (8 - 2x + z) + 2z &= 10,\\[-4pt]
\multicolumn{2}{c}{\rule{6cm}{0.3pt}}\\[-2pt]
x - 16 + 4x - 2z + 3z &= -4,\\
3x + 8 - 2x + z + 2z &= 10,\\[-4pt]
\multicolumn{2}{c}{\rule{6cm}{0.3pt}}\\[-2pt]
5x + z &= 12,\\
x + 3z &= 2.
\end{align*}

Máme tedy soustavu dvou rovnic o dvou neznámých:
\[
\begin{cases}
5x + z = 12,\\
x + 3z = 2.
\end{cases}
\]

\textit{Upravíme druhou rovnici a přičteme:}
\begin{align*}
5x + z &= 12,\\
x + 3z &= 2 \quad | \cdot(-5),\\[-4pt]
\multicolumn{2}{c}{\rule{6cm}{0.3pt}}\\[-2pt]
5x + z - 5x - 15z &= 12 - 10,\\
-14z &= 2 \quad | :(-14),\\
z &= -\tfrac{1}{7}.
\end{align*}

Dosadíme \(z = -\tfrac{1}{7}\) do \(x + 3z = 2\):
\begin{align*}
x + 3z &= 2,\\
x + 3\left(-\tfrac{1}{7}\right) &= 2,\\
x - \tfrac{3}{7} &= 2 \quad | +\tfrac{3}{7},\\
x &= \tfrac{17}{7}.
\end{align*}

Dosadíme \(x, z\) do \(y = 8 - 2x + z\):
\begin{align*}
y &= 8 - 2x + z,\\
y &= 8 - 2\cdot\tfrac{17}{7} - \tfrac{1}{7},\\
y &= \tfrac{56 - 34 - 1}{7},\\
y &= \tfrac{21}{7} = 3.
\end{align*}

\textbf{Řešení soustavy:} \quad \(x = \tfrac{17}{7},\; y = 3,\; z = -\tfrac{1}{7}\)

\medskip
\textbf{Množina řešení:} \quad \(K = \left\{\left[\tfrac{17}{7};\,3;\,-\tfrac{1}{7}\right]\right\}\)
\end{examplebox}







% -----------------------------------------------------------------
% ---------------------- P Ř Í K L A D Y --------------------------
% -----------------------------------------------------------------

\exercisepart
\setcounter{section}{1} % případně uprav dle předchozího číslování


% =========================
% ÚLOHA 1
% =========================
\begin{exercise}[Lineární rovnice]
  Řešte následující jednoduché lineární rovnice s neznámou \(x \in \mathbb{R}\) a proveďte zkoušku správnosti řešení.
  \begin{tasks}(3)
    \task $2x - 7 = 5$
    \task $3x + 6 = 8$
    \task $\pi x + 3= 4$
    \task \(\dfrac{x}{2} + 6 = 12\)
    \task \(0{,}83\,x + 32 = 160\)
    \task $1-\sqrt{3}x=\sqrt{3}$
    \task \(5 - x\sqrt{3} = 2\)
    \task \(x\sqrt{5} - \sqrt{5} = 5\)
    \task $2x - \sqrt{5} = \sqrt{3}$
  \end{tasks}
\end{exercise}


% =========================
% ÚLOHA 2
% =========================
\begin{exercise}[Lineární rovnice se zlomky]
  Řešte následující rovnice obsahující zlomky a proveďte zkoušku.
  \begin{tasks}(2)
    \task \(\dfrac{3y - 1}{5} = \dfrac{1}{10}\)
    \task \(2 + \dfrac{y - 1}{2} - y = 1\)
    \task \(\dfrac{4x - 5}{2} = x - \left(\dfrac{5}{2} - x\right)\)
    \task \(\dfrac{4x - 5}{2} = x - (1 - x)\)
    \task \(\dfrac{3 - 2x}{4} + \dfrac{x}{2} = 1\)
    \task \(\dfrac{5y + 1}{3} - \dfrac{y - 2}{6} = 2\)
    \task \(\dfrac{7 - 3t}{5} = \dfrac{t + 8}{10}\)
    \task \(1 - \dfrac{2x - 3}{4} = \dfrac{x}{2}\)
    \task \(\dfrac{3(x - 1)}{8} + \dfrac{1}{2} = \dfrac{x}{4}\)
    \task $\dfrac{4x - 7}{2} - \dfrac{x - 4}{6} = 2x - 3 $
    %\task $\dfrac{x+1}{x-1} + \dfrac{2}{x+2} - 1 = \dfrac{6}{x^2 + x - 2} $
  \end{tasks}
\end{exercise}


% =========================
% ÚLOHA 3
% =========================
\begin{exercise}[Lineární rovnice s neznámou na obou stranách]
  Vyřešte následující rovnice, ve kterých se neznámá nachází na obou stranách rovnosti.  
  Každé řešení ověřte zkouškou dosazením do původní rovnice.
  \begin{tasks}(2)
    \task \(2x - 3 = 5x + 1\)
    \task \(\pi x + 4 = x - \pi\)
    \task \(x\sqrt{2} + 7 = x\sqrt{3} - 2\)
    \task \(10x + 5 = 20x - 5\)
    \task \(4 - \sqrt{2}\,x = 2x + 1\)
    \task \((3 - \pi)x + 2 = (1 + \pi)x - 5\)
    \task \(5(2x - 1) = 3x + 7\)
    \task \(\sqrt{3}\,x - 2 = \sqrt{5}\,x + 1\)
    \task \(7x + 4 = 3(2x + 1) + 1\)
    \task* \(3[2(3x - 6) - 2(4x - 5) + 1] = 6[3 - 8(x - 3)]\)
  \end{tasks}
\end{exercise}


% =========================
% ÚLOHA 4
% =========================
\begin{exercise}[Lineární rovnice]
  Řešte dané rovnice.
  \begin{tasks}(2)
    \task \(3(2x + 5) + 15 = x + 85\)
    \task \(2(2x + 1) - (1 - x) = 5(x + 3)\)
    \task \((x - 2)^2 + (x + 1)^2 = (2x + 3)(x - 1)\)
    \task \(4(x + 3)^2 - (2x + 1)^2 = 4(5x + 8) + 3\)
    \task \((3x + 1)^2 + (4x + 1)^2 = (5x + 1)^2 + 1\)
  \end{tasks}
\end{exercise}


% =========================
% ÚLOHA 5
% =========================
\begin{exercise}[Složené rovnice a identity]
  Řešte následující rovnice obsahující složené výrazy a závorky různých řádů.  
  Použijte vhodné úpravy a proveďte zkoušku správnosti řešení.
  \begin{tasks}(1)
    \task \((x-3)(x+2) = (x-2)(x-1)\)
    \task \(5\{ 5[ 5(5x-4)-4] -4\}  = 5\)
    \task \(3\bigl[\,2(3x - 6) - 2(4x - 5) + 1\,\bigr] = 6\bigl[\,3 - 8(x - 3)\,\bigr]\)
    \task \(x - (2x - 0{,}2)^2 = 5{,}8 - (2x + 0{,}1)(2x - 0{,}1)\)
    \task \((x-1)^3 + (x-2)^3 + (x-3)^3 = 3(x-1)(x-2)(x-3)\)
    \task \(\dfrac{x}{2} - \dfrac{s - \tfrac{s}{2}}{2}
            - \dfrac{\left(x - \tfrac{s}{2}\right) - \tfrac12\left(2 + \tfrac{x}{2}\right)}{2}
            = \tfrac12\!\left(x - \tfrac{x}{2}\right)\!\cdot \tfrac12\)
  \end{tasks}
\end{exercise}


% =========================
% ÚLOHA 6
% =========================
\begin{exercise}[Rovnice MIX]
  Vyřešte následující rovnice.
  \begin{tasks}
    \task \((x + 1)^3 - (x - 1)^3 = 6(x + 2)(x - 1) + 9(x + 1) - 9(x - 1)\)
    \task \(\dfrac{3 - 2x}{2} - \left(\dfrac{3 + x}{3} - \dfrac{2x - 3}{4}\right) + 2
    = x - \left(\dfrac{x + 1}{6} - \dfrac{2 - x}{12}\right)\)
    \task \(\dfrac{x + \dfrac{x + \dfrac{2x + 3}{2}}{3}}{3} = \dfrac{x - \dfrac{x - \dfrac{x - 2}{2}}{3}}{4}\) % pro vizuálně čistou strukturu
    \task \(\dfrac{9n - 0{,}7}{4} - \dfrac{5n - 1{,}5}{7}  = \dfrac{7n - 1{,}1}{3} - \dfrac{5(0{,}4 - 2n)}{6}\)
  \end{tasks}
\end{exercise}


% =========================
% ÚLOHA 7
% =========================
\begin{exercise}[Rovnice se zlomky – obtížnější]
  Vyřešte následující rovnice a výsledné řešení ověřte zkouškou.
  \begin{tasks}
    % dlouhá rovnice přes celou šířku
    \task \(\dfrac{4-x}{2}
              - \left(\dfrac{2+7x}{8} - \dfrac{8-x}{6}\right)
              + \dfrac{x+2}{4}
              - \dfrac{8-x}{3}
              + x - 1 = 0\)

    \task \(\dfrac{y-4}{5} - \dfrac{7(3y-2)}{5}
           = \dfrac{9}{13}\,(9-2y) - (3y-5)\)

    \task \(\dfrac{3(z+1)}{2} - \left(\dfrac{z+9}{4} - 1\right)
           = \dfrac{5z+1}{7} - \left(\dfrac{3z+1}{2} - 4\right)\)

    % podobný nový příklad
    \task \(\dfrac{2(t-3)}{5}
           - \left(\dfrac{t+4}{3} - \dfrac{2-t}{6}\right)
           + \dfrac{3t-1}{4}
           = \dfrac{t}{2} - \left(\dfrac{5t-7}{10} - 1\right)\)
  \end{tasks}
\end{exercise}

\newpage
% =========================
% ÚLOHA 8
% =========================
\begin{exercise}[Rovnice s mocninami a rozvinutím závorek]
  Řešte následující rovnice, ve kterých vystupují kvadráty a kubické členy a proveďte následné zjednodušení výrazů.
  \begin{tasks}(1)
    \task \((3u - 1)^2 - 5(2u + 1)^2 + 2(6u - 3)(u + 1) = (u - 1)^2 + 3(2u - 1)\)
    \task \((v + 1)^3 - (v - 1)^3 = 6(v + 2)(v - 1) - 2(v + 1) + 2(v - 1)\)
    \task \((2x - 3)^2 - 4(x + 1)(x - 2) = (x - 3)^2 - (x + 2)^2 + 5x\)
    \task \((4y + 1)^2 - 2(3y - 2)^2 + 5(2y + 1) = (y - 1)(5y + 3)\)
    \task \((p - 2)^3 + (p + 1)^3 = 3(p - 1)(p + 2)(p - 2) + 9p\)
    \task \((t + 2)^2 - (t - 3)^2 = 5(2t - 1) - (t + 4)(t - 2)\)
  \end{tasks}
\end{exercise}


% =========================
% ÚLOHA 9
% =========================
\begin{exercise}[Součinové rovnice]
  Vyřešte následující součinové rovnice.
  \begin{tasks}(2)
    \task \((x+1)(x-3)(x+\pi)=0\)
    \task \(x(3x+1)(x+\sqrt{2})(4x-\pi)=0\)
    \task \((3x+2)(x\sqrt{2}+1)(\pi^{2}x+\pi)=0\)
    \task $\bigl(x\sqrt{2}-x-1\bigr)\,(\pi x+\sqrt{2})\,(x^{2}+1)=0$
    \task \((2x-5)(x+1)(x-\pi)=0\)
    \task \((x-4)(\sqrt{3}\,x+2)(x^{2}+4)=0\)
    \task $\bigl(5x+\pi\bigr)\bigl(x-\sqrt{5}\bigr)(2x+1)=0$
    \task \(x(x-1)(x+1)(\pi x-2)=0\)
    \task \((3x-2)\bigl(x+\sqrt{7}\bigr)\,(\pi^{2}x-1)=0\)
    \task $9 - x^2 = 3 - x$
  \end{tasks}
\end{exercise}


\setcounter{section}{2}
% =========================
% ÚLOHA 10
% =========================
\begin{exercise}[Rovnice s absolutní hodnotou]
  Řešte následující rovnice s absolutní hodnotou neznámé \(x \in \mathbb{R}\). 
  Každé řešení ověřte zkouškou.
  \begin{tasks}(3)
    \task \(|x| = 7\)
    \task \(|x| = \tfrac{2}{3}\)
    \task \(|x - 1| = 3\)
    \task \(|x + 4| = 1\)
    \task \(|x + \pi| = 6\)
    \task \(|2x - 3| = 6\)
    \task \(|6 - x| = 2\)
    \task \(|\,x - \sqrt{3} + 1\,| = 2 - \sqrt{2}\)
    \task \(|x - 5\sqrt{11}| = 0\)
    \task \(|4x - 7| = -1\)
    \task $|x - \sqrt{3}| = 2 + 5\sqrt{3}$
    \task $\dfrac{|x+3|}{2}=1$
  \end{tasks}
\end{exercise}


% =========================
% ÚLOHA 11
% =========================
\begin{exercise}[Rovnice s absolutní hodnotou a lineárním výrazem]
  Vyřešte následující rovnice.
  \begin{tasks}(2)
    \task \(|x - 7| = x - 7\)
    \task \(|x - 2| = 2 - x\)
    \task \(|8 - 5x| = 5x - 8\)
    \task \(|2x - 3| = x\)
    \task $|x + 4| = 3x$
    \task \(|2x - 5| = 1 - 3x\)
  \end{tasks}
\end{exercise}


% =========================
% ÚLOHA 12
% =========================
\begin{exercise}[Rovnice s více absolutními hodnotami]
  Řešte následující rovnice, které obsahují dvě nebo více absolutních hodnot a kontrolujte správnost nalezených řešení.
  \begin{tasks}(2)
    \task \(|4 - x| - |2x + 3| = 7\)
    \task $|2x - 1| - |3 - x| = x$
    \task \(|2x - 4| - |x + 3| = 2 - |x - 5|\)
    \task \(|x - 1| + 3|2 - x| = x - |1 - x|\)
    \task \(|x - 4| + |2x - 1| = |x| + 3\)
    \task \(|x + 2| = 4|x - 3|\)
  \end{tasks}
\end{exercise}

\setcounter{section}{3}
% =========================
% ÚLOHA 13
% =========================
\begin{exercise}[Základní nerovnice]
  Vyřešte následující jednoduché nerovnice v oboru reálných čísel \(\mathbb{R}\).  
  Určete množinu všech řešení a ověřte, zda mají smysl.
  \begin{tasks}(3)
    \task \(0\cdot x \le 0\)
    \task \(0\cdot x < 0\)
    \task \(0\cdot x \le -3\)
    \task \(3x < 0\)
    \task \(-5x \le 0\)
    \task $2 \ge 0$
    \task $-3x < 0$
    \task $2x \le 2$
    \task $0x \ge 0$
  \end{tasks}
\end{exercise}


% =========================
% ÚLOHA 14
% =========================
\begin{exercise}[Lineární nerovnice s různými neznámými]
  Řešte následující lineární nerovnice obsahující jednu neznámou.
  \begin{tasks}(3)
    \task \(7 - 2x > 0\)
    \task \(y + 1 \ge y\)
    \task \(x - 7 < 8 - x\)
    \task \(2z + 1 > 1 - z\)
    \task \(3x - 7 \le 5x - 13\)
    \task \(x + 3 \le 2x - 7\)
    \task \(2y + 5 < 2y + 3\)
    \task \(5(x + 1) < 2x - 4\)
    \task $4x + 1 > 2x + 3$
    \task \(x + 3 \ge 7 + \sqrt{2}\,x\)
    \task \(\sqrt{3}\,x + 3 \ge 2 - x\)
    \task \(x\sqrt{2} + 1 > \sqrt{3} + x\)
    \task \(2x - \sqrt{2} < x\sqrt{2} - 2\)
    \task \(2\pi z > 1\)
    \task \(3x - 1 \le 5 + \pi x\)
  \end{tasks}
\end{exercise}


% =========================
% ÚLOHA 15
% =========================
\begin{exercise}[Lineární nerovnice se zlomky]
  Řešte následující nerovnice obsahující zlomky.  
  \begin{tasks}(2)
    \task \(\dfrac{4x}{3} \le \dfrac{2}{3} + x\)
    \task $\dfrac{2+3x}{18} \le 0$
    \task $\dfrac{x+2}{3} > \dfrac{x-2}{3}$
    \task $\dfrac{x+2}{3} \le \dfrac{x-2}{3}$
    \task \(2 - \dfrac{x+2}{3} > x - \dfrac{x+3}{3}\)
    \task \(6x + 1 > 2(x - 5) - 1\)
    \task \(\dfrac{4x - 7}{2} - \dfrac{x - 4}{3} \ge 2x - 3\)
    \task \(\dfrac{x+1}{4} - \dfrac{x-3}{3} < \dfrac{x+4}{6} + \dfrac{7 - 3x}{12}\)
    \task $\dfrac{7}{6} - \dfrac{x}{2} > \dfrac{3}{7}x + \dfrac{39}{14}-x$
    \task $\dfrac{2x}{3} - 3 > \dfrac{16x}{21} - \dfrac{13}{3} - \dfrac{2x}{15}$
  \end{tasks}
\end{exercise}


% =========================
% ÚLOHA 16
% =========================
\begin{exercise}[Kvadratické a mocninné nerovnice]
  Vyřešte následující nerovnice obsahující kvadráty či vyšší mocniny.
  \begin{tasks}(2)
    \task \((4 - t)t > 24 - t^{2}\)
    \task \((3n - 5)^{2} + (4n - 3)^{2} \ge (5n - 4)^{2}\)
    \task \((u - 3)^{2} + 3(u - 1) \le (u + 3)(u - 3)\)
    \task \((x - 1)^{2} - (x + 1)^{2} < 8\)
    \task \((x + 1)^{2} - (x - 1)^{2} + 2 \ge 2(x + 1)\)
    \task \((x - 2)^{2} \ge (x + 1)(x - 5)\)
  \end{tasks}
\end{exercise}


% =========================
% ÚLOHA 17
% =========================
\begin{exercise}[Součinové nerovnice – intervalová metoda]
  Řešte následující součinové nerovnice pomocí metody intervalů.  
  \begin{tasks}(2)
    \task \((x-2)(x+1) > 0\)
    \task \((x-2)(x+3) \ge 0\)
    \task \((y+3)\!\left(y-\tfrac{1}{2}\right) \ge 0\)
    \task \(\left(z+\tfrac{1}{2}\right)\!\left(z-\tfrac{1}{2}\right) \le 0\)
    \task \((1-x)(x+\sqrt{2}) > 0\)
    \task \((3-5y)(3+5y) \le 0\)
    \task \(\left(2z-\tfrac{1}{3}\right)\!\left(z+\tfrac{2}{3}\right) < 0\)
    \task \((2u-5)\bigl(3u-\sqrt{2}\bigr) \le 0\)
  \end{tasks}
\end{exercise}


% =========================
% ÚLOHA 18
% =========================
\begin{exercise}[Součinové nerovnice vyššího řádu]
  Vyřešte následující součinové nerovnice s více činiteli nebo s vyššími mocninami.
  \begin{tasks}(2)
    \task \((x-1)(x-2)(x-3) > 0\)
    \task \((2x+1)2x(2x-1) \le 0\)
    \task \((2-z)(3-z)(4-z)(5-z) < 0\)
    \task \((3-z)(2-z)(1-z)\,z \ge 0\)
    \task \((4x^{2}-3)\,(4-5x)\,(x-2) < 0\)
    \task \((3-x)^{2}(x-1)\,(\sqrt{2}-x) \ge 0\)
    \task \((4-7x)(x+1)(10x-7)(\pi-x) \le 0\)
  \end{tasks}
\end{exercise}


\setcounter{section}{4}
% =========================
% ÚLOHA 19
% =========================
\begin{exercise}[Lineární nerovnice s absolutní hodnotou]
  Řešte následující jednoduché nerovnice s absolutní hodnotou pro \(x \in \mathbb{R}\).  
  \begin{tasks}(3)
    \task \(|x| < 5\)
    \task \(|x| \ge 3\)
    \task \(|x - 2| \le 4\)
    \task \(|x + 3| > 6\)
    \task \(|x - \pi| < 2\)
    \task \(|x - \sqrt{3}| \ge 1\)
    \task \(|2x - 5| < 3\)
    \task \(|3x + 2| \ge 4\)
    \task \(|x - 1| + 2 > 5\)
  \end{tasks}
\end{exercise}


% =========================
% ÚLOHA 20
% =========================
\begin{exercise}[Složenější nerovnice s absolutní hodnotou]
  Vyřešte následující nerovnice, které obsahují více absolutních hodnot.  
  \begin{tasks}(2)
    \task \(|2x - 3| \le |x + 1|\)
    \task \(|x - 4| > |2x + 1|\)
    \task \(|x - 1| + |x + 2| \le 5\)
    \task \(|2x - 3| - |x + 1| > 1\)
    \task \(|x + 1| + 1 < 2|x - 1|\)
    \task \(|3x - 2| \ge 2|x + 1|\)
  \end{tasks}
\end{exercise}


% =========================
% ÚLOHA 21
% =========================
\begin{exercise}[Nerovnice s absolutní hodnotou]
  Vyřešte následující nerovnice s absolutní hodnotou převodem na ekvivalentní soustavy lineárních nerovnic.
  \begin{tasks}(2)
    \task \(|x - 2| < 3x - 1\)
    \task \(|2x + 1| > 3 - x\)
    \task \(|x - 5| \le 2x + 1\)
    \task \(|3x - 2| \ge 2x - 1\)
    \task \(|x + 4| + 1 \le 3x\)
    \task \(|2x - 1| > |x| + 2\)
  \end{tasks}
\end{exercise}


\setcounter{section}{5}
% =========================
% ÚLOHA 22
% =========================
\begin{exercise}[Racionální rovnice]
  Řešte následující rovnice, určete definiční obor a proveďte zkoušku.
  \begin{tasks}(3)
    \task \(\dfrac{2x + 5}{3x - 6} = 0\)
    \task \(\dfrac{3x - 2}{x + 5} = 0\)
    \task \(\dfrac{7x + 3}{x + \sqrt{2}} = 0\)
    \task \(\dfrac{4x - 9}{2x - 3} = 0\)
    \task \(\dfrac{x - 5\sqrt{3}}{x - 1} = 0\)
    \task \(\dfrac{5x - 4}{7x + 2} = 0\)
    \task \(\dfrac{2x + 7}{x - \pi} = 0\)
    \task \(\dfrac{9 - 3x}{2x + 1} = 0\)
    \task \(\dfrac{4x + 1}{x - \tfrac{1}{2}} = 0\)
  \end{tasks}
\end{exercise}



% =========================
% ÚLOHA 23
% =========================
\begin{exercise}[Racionální rovnice]
  Řešte rovnice, ve kterých je zlomek roven dané konstantě.  
  Určete definiční obor a proveďte zkoušku správnosti řešení.
  \begin{tasks}(3)
    \task \(\dfrac{3x + 6}{x + 2} = 3\)
    \task \(\dfrac{3x + 6}{x + 2} = 5\)
    \task \(\dfrac{x - 3}{x - 4} = \tfrac{1}{2}\)
    \task \(\dfrac{x + 2}{3x - 1} = \tfrac{1}{3}\)
    \task \(\dfrac{x - 3}{x - 4} = 1\)
    \task \(\dfrac{x - \sqrt{5}}{x + \sqrt{5}} = \sqrt{5}\)
    \task \(\dfrac{2x - 7}{x + 1} = -3\)
    \task \(\dfrac{4x + 1}{2x - 5} = 2\)
    \task \(\dfrac{5x - 3}{x - 4} = \pi\)
  \end{tasks}
\end{exercise}



% =========================
% ÚLOHA 24
% =========================
\begin{exercise}[Racionální rovnice – rovnost dvou podílů]
  Vyřešte rovnice, ve kterých se rovnají dva racionální výrazy.  
  Určete podmínky jmenovatelů a proveďte křížové násobení.
  \begin{tasks}(2)
    \task \(\dfrac{x - 1}{x - 3} = \dfrac{x + 3}{x - 6}\)
    \task \(\dfrac{(3 - x)(2 + x)}{(4 + x)(x - 3)} = \dfrac{2 + x}{x - 3}\)
    \task \(\dfrac{\tfrac{1}{3}y + \tfrac{3}{4}}{\tfrac{1}{3}y - \tfrac{3}{4}} = \dfrac{y + 3}{y - 4}\)
    \task \(\dfrac{2x + 1}{x - 2} = \dfrac{3x - 4}{2x + 5}\)
    \task \(\dfrac{x - 5}{x + 2} = \dfrac{2x + 3}{x - 1}\)
    \task \(\dfrac{3x + 4}{x - \pi} = \dfrac{x - 2}{x + \pi}\)
    \task \(\dfrac{x + \sqrt{2}}{2x - 3} = \dfrac{3x - 1}{x + \sqrt{2}}\)
    \task \(\dfrac{2x - 5}{x + 1} = \dfrac{3x + 4}{2x - 3}\)
    \task \(\dfrac{n + \sqrt{2}}{n - \sqrt{2}} = \dfrac{49 - 14\sqrt{2}}{47}\)
    \task \(\dfrac{n + \sqrt{2}}{n - \sqrt{2}} = \dfrac{51 + 14\sqrt{2}}{47}\)
  \end{tasks}
\end{exercise}



% =========================
% ÚLOHA 25
% =========================
\begin{exercise}[Racionální rovnice – součty a rozdíly podílů]
  Řešte rovnice, které obsahují součty nebo rozdíly racionálních výrazů.  
  Stanovte podmínky definičního oboru a proveďte zkoušku.  
  Využijte například rozklad \(x^{2} + x - 2 = (x + 2)(x - 1)\).
  \begin{tasks}(2)
    \task \(\dfrac{x - 1}{x - 2} + \dfrac{x + 1}{x - 4} = 2\)
    \task $\dfrac{5}{x+1} - 1 = \dfrac{8-x}{x-1}$
    \task \(\dfrac{2x + 3}{x - 1} - \dfrac{x - 5}{x + 2} = 4\)
    \task $5 + \dfrac{3}{3x-12} = \dfrac{5-x}{x-4}$
    \task $\dfrac{3}{x+1} - 1 = \dfrac{6-3x}{x-1}$
    \task $\dfrac{4x-1}{x-2} + 1 = \dfrac{3-5x}{2-x}$
    \task \(\dfrac{3x - 2}{x + 4} + \dfrac{x + 6}{x - 4} = 5\)
    \task \(\dfrac{x + 2}{x - 2} - \dfrac{x - 2}{x + 2} = 1\)
    \task $\dfrac{x-4}{2(x-1)} + \dfrac{x+4}{2(x+1)} = 1$
    \task \(\dfrac{2}{x} + \dfrac{3}{x - 1} = \dfrac{5}{x^{2} - x}\)
    \task \(\dfrac{x + 1}{x + 3} + \dfrac{x - 1}{x - 3} = \dfrac{2x}{x^{2} - 9}\)
    \task \(\dfrac{1}{x - 3} + \dfrac{2}{x + 1} = \dfrac{3}{x^{2} - 2x - 3}\)  
    \task \(\dfrac{x + 1}{x - 1} + \dfrac{2}{x + 2} - 1 = \dfrac{6}{x^{2} + x - 2}\)  
    \task $\frac{-1 + k + \dfrac{2 - x}{3}}{-1 - k + \dfrac{2 + k}{4}} = \dfrac{8}{11}$
  \end{tasks}
\end{exercise}


% =========================
% ÚLOHA 26
% =========================
\begin{exercise}[BONUS: Racionální rovnice]
  Řešte rovnici a proveďte zkoušku.
  \begin{tasks}(2)
    \task \(\dfrac{9}{x - \dfrac{3}{x - \dfrac{9}{x - 3}}} = 3\)
    \task \(\dfrac{5}{5 - \dfrac{30}{6 - \dfrac{5}{6 - x}}} = x\)
    \task \(\dfrac{8}{2 - \dfrac{16}{x + 1 - \dfrac{4}{x}}} = x\)
    \task \(\dfrac{4}{3 - \dfrac{12}{2 - \dfrac{6}{x - 1}}} = x\)
  \end{tasks}
\end{exercise}


\setcounter{section}{6}
% =========================
% ÚLOHA 27
% =========================
\begin{exercise}[Lineární rovnice s více neznámými]
  Řešte následující lineární rovnice s více neznámými.
  \begin{tasks}(2)
    \task \(2x + 3y = 12\)
    \task $x - \frac{1}{2}x = 4$
    \task $-3x - 2y = -5$
    \task $0x + 4y = 7$
    \task $-x + 0y = 3$
    \task $0x - 2y = 0$
    \task $0x + 0y = -4$
    \task \(2x + y - 2 = x + 3y + 3\)
    \task \(3(x+y) - x + 2 = \sqrt{2}\,x + y\)
    \task \(-\tfrac{3}{4}x - y = \tfrac{3}{2} - y\)
    \task \(4x - y + 10 = y + 10\)
    \task \(\dfrac{x+1}{\,y-3\,} = 2\)
    \task \(2x - 3y + z - 5 = 0\)
    \task \(x - 2y + 3z + 3 = 2x + 3y + 3z\)
  \end{tasks}
\end{exercise}


\setcounter{section}{7}

% =========================
% ÚLOHA 28
% =========================
\begin{exercise}[Soustavy lineárních rovnic – základní]
Řešte následující soustavy rovnic v reálných číslech. Zapište řešení (pokud existuje) jako uspořádanou dvojici/trojici, případně popište množinu řešení.
\begin{tasks}(3)

\task $\begin{cases}
-x + 5y = -12\\
x - 4y = 15
\end{cases}$

\task $\begin{cases}
3u + 2v = -4\\
2u - 3v = 19
\end{cases}$

\task $\begin{cases}
7x - 3y = 15\\
5y +6y = 27
\end{cases}$

\task $\begin{cases}
3x + 2y = 20\\
2x + 3y = 20
\end{cases}$

\task $\begin{cases}
3u + 2v = -4\\
2u - 3v = 19
\end{cases}$

\task $\begin{cases}
3u + 2v = -4\\
2u - 3v = 19
\end{cases}$

\task $\begin{cases}
3u + 2v = -4\\
2u - 3v = 19
\end{cases}$

\task $\begin{cases}
3u + 2v = -4\\
2u - 3v = 19
\end{cases}$

\task $\begin{cases}
u - 2v = \sqrt{2}\\
2u + v = 3\sqrt{2}
\end{cases}$

\task $\begin{cases}
2x - 3y = 1\\
-x + \tfrac{3}{2}y = -\tfrac{1}{2}
\end{cases}$

\task $\begin{cases}
y + \sqrt{3}\,z = 1\\
2y + 2\sqrt{3}\,z = 3
\end{cases}$

\task $\begin{cases}
0\cdot x + 3y = 1\\
-x - y = 0
\end{cases}$

\task $\begin{cases}
\sqrt{2}\,x + 2y = -2\\
3x + 0\cdot y = 4
\end{cases}$

\task $\begin{cases}
u = 4 - 2v\\
v = 3 + u
\end{cases}$

\task $\begin{cases}
4t = 5 + 2s\\
6 + 2x = 4t
\end{cases}$

\end{tasks}
\end{exercise}



\newpage
% =========================
% ÚLOHA 29
% =========================
\begin{exercise}[Soustavy lineárních rovnic – obtížnější a směs typů]
Řešte soustavy.
\begin{tasks}(1)

\task $\begin{cases}
3(x+4) - 2(x+y) = 4\\
5(x-y) + 6(x-2) - 2 = 0
\end{cases}$

\task $\begin{cases}
(x+4)(s-2) - (x-5)(s+4) = 0\\
(x+6)(s-1) - (x-1)(s+2) = 0
\end{cases}$

\task $\begin{cases}
(x+4)(y-2) = (x-2)(y+13) \\
(x-1)(y-3) = (x+2)(y-5)
\end{cases}$

\task $\begin{cases}
\tfrac{1}{3}(2x-3y) - \tfrac{1}{4}(x-2y+3) = 3\\
\tfrac{3}{4}(x+1) + \tfrac{1}{3}(4x-2y-3) = 6 + y
\end{cases}$

\task $\begin{cases}
(x-1)(y+4) = (x+3)(y-1)\\
(x-2)(y+2) = x(y-2)
\end{cases}$

\task $\begin{cases}
(x-y)^2 - (x+y)^2 = 2x(1-2y) - y\\
(x-y)(x+y) - 2y(x-1) = (x-y)^2 + 4x - 2y^2
\end{cases}$

\task $\begin{cases}
\sqrt{2}\,x - \sqrt{3}\,y = 1\\
(2\sqrt{2}-\sqrt{3})x + (\sqrt{2}-2\sqrt{3})y = 0
\end{cases}$

\task $\begin{cases}
\tfrac{x}{2} - \tfrac{y}{3} = 1\\
\tfrac{3x}{4} + \tfrac{y}{6} = \tfrac{5}{2}
\end{cases}$

\task $\begin{cases}
(x+1)(y-2) - (x-1)(y+2) = 7\\
(2x-3)(y+1) + (x+4)(y-2) = 1
\end{cases}$

\end{tasks}
\end{exercise}


% =========================
% ÚLOHA 31
% =========================
\begin{exercise}[Soustavy lineárních rovnic – soustavy o třech neznámých]
Řešte následující soustavy rovnic. Uveďte, zda mají právě jedno řešení, žádné řešení, nebo nekonečně mnoho řešení (a případně parametrický tvar).
\begin{tasks}(3)

\task $\begin{cases}
4x + y + 2z = 2\\
2y - z = 1\\
z = 7
\end{cases}$

\task $\begin{cases}
x + y + z = 1\\
x + 2y + 2z = 2\\
2x + 2y + z = 3
\end{cases}$

\task $\begin{cases}
x - y + z = 2\\
x + 2y - 2z = 3\\
6x - 3y + 3z = 14
\end{cases}$

\task $\begin{cases}
x - y + 2z = 2\\
x + 2y - 2z = 3\\
6x - 3y + 8z = 13
\end{cases}$

\task $\begin{cases}
3x + 2y + 3z = 10\\
5x - y + 4z = 2\\
2x - 3y + 5z = -3
\end{cases}$

\task $\begin{cases}
-2x + y + 4z = -2\\
x - 2y - 2z = 0\\
x + y - 2z = 2
\end{cases}$

\task $\begin{cases}
5x + 5y + 4z = 2\\
3x + 4y - 3z = 1\\
-2x - y - 4z = 0
\end{cases}$

\task $\begin{cases}
2x + 3y = 2\\
3x + 2z = 1\\
3y + 4z = 2
\end{cases}$

\task $\begin{cases}
x + y + 2z = -1\\
2x - y + 2z = -4\\
4x + y + 4z = -2
\end{cases}$

\task $\begin{cases}
2x + y - z = 0\\
4x + 2y + z = 0\\
x - y + 3z = 0
\end{cases}$

\task $\begin{cases}
2x - y = 6\\
y + 4z = 8\\
x - z = 1
\end{cases}$

\task $\begin{cases}
(x-2z)+(y+z)=3\\
2(x-y)-(y-3z)=4\\
x+y+z=2
\end{cases}$

\task $\begin{cases}
x-2y+z=1\\
-x+3y+2z=5\\
2x-y+5z=5
\end{cases}$


\task $\begin{cases}
-2y+z=1\\
2x+y-2z=5\\
x+2y-4z=-14
\end{cases}$

\task $\begin{cases}
x+y-2z=-2\sqrt{2}\\
2x+2y+3z=3\sqrt{2}\\
5x+5y+4z=4\sqrt{2}
\end{cases}$

\end{tasks}
\end{exercise}


% =========================
% ÚLOHA 31
% =========================
\begin{exercise}[Řešte soustavy rovnic]
\begin{tasks}(2)

\task $\begin{cases}
\dfrac{x+1}{y+1} = 2\\[6pt]
\dfrac{y+2}{z+1} = 4\\[6pt]
\dfrac{z+3}{x+1} = \frac{1}{2}\\[6pt]
\end{cases}$

\task $\begin{cases}
\dfrac{3x+y}{z+1} = 2\\[6pt]
\dfrac{3y+z}{x+1} = 2\\[6pt]
\dfrac{3z+x}{y+1} = 2
\end{cases}$

\end{tasks}
\end{exercise}


% =========================
% ÚLOHA 32
% =========================
\begin{exercise}[Řešte soustavy rovnic se čtyřmi neznámými \(x, y, z, u\)]
\begin{tasks}(2)

\task $\begin{cases}
x + y + z + u = 10\\
x + y - z - u = -4\\
x - y - z + u = 0\\
-x + y + z + u = 8
\end{cases}$

\task $\begin{cases}
x - 2z + u = 3\\
2x - y + z = 5\\
x - y - 2u = 1\\
2y + z - u = -3
\end{cases}$

\end{tasks}
\end{exercise}


\end{document}
